\section{Source front ends and metacircular compilation}
\label{sec:front-end}

This section closes the boundary left abstract in
Equation~\eqref{eq:compilation-unit}.  It has two separate purposes.  First,
it defines the relation from source text through PEG recognition, surface AST
construction, stable source identities, and elaboration to the annotated core
used by the preceding rules.  Second, it makes the implementations of those
stages ordinary Newspeak objects.  The latter introduces no new reflective
command: changing a parser, elaborator, or compiler is an ordinary change to
its methods or object graph under Section~\ref{sec:reflection}.

\subsection{Source text and PEG recognition}
\label{sec:peg-recognition}

Let $u\in\mathit{SourceText}=\mathit{UnicodeScalar}^{*}$, let
$0\le i\le |u|$ be a source offset, let $N\in\mathit{Nonterminal}$, and let
$X\subseteq\mathit{UnicodeScalar}$.  Character ranges, optional expressions,
positive closure, and the Specification's prefix \texttt{\textasciitilde}
notation are abbreviations for the following PEG core:
\begin{equation}\label{eq:peg-expression-domain}
 r::=\epsilon\mid\kw{char}(c)\mid\kw{set}(X)\mid\kw{any}
   \mid N\mid r_1r_2\mid r_1/r_2\mid r^{*}
   \mid {!r}\mid {\&r}\mid\kw{capture}(n,r).
\end{equation}
A grammar is $G=\langle N_0,\rho_G\rangle$, where $N_0$ is its start symbol
and $\rho_G:\mathit{Nonterminal}\rightharpoonup\mathit{PEGExpr}$ is finite.
A concrete tree is
\begin{equation}\label{eq:concrete-tree-domain}
 t::=\kw{leaf}\langle i,j\rangle
   \mid\kw{node}\langle n,i,j,\bar t\rangle,
 \qquad \mathsf{span}(t)=[i,j).
\end{equation}
Uncaptured recognition returns an ordered forest.  Write
$\mathsf{ok}(j,\bar t)$ for success and $\mathsf{fail}(i)$ for failure at
offset $i$.  The deterministic function
$\mathsf{peg}_{G}(r,u,i)$ is defined by the following exhaustive equations.
The failure offset is diagnostic only; ordered choice always restarts its
right alternative at the original offset.
\begin{equation}\label{eq:peg-atomic-recognition}
\begin{aligned}
 \mathsf{peg}_{G}(\epsilon,u,i)&=\mathsf{ok}(i,\epsilon),\\
 \mathsf{peg}_{G}(\kw{char}(c),u,i)&=
  \begin{cases}\mathsf{ok}(i+1,\langle\kw{leaf}\langle i,i+1\rangle\rangle)
       &i<|u|\land u_i=c,\\
    \mathsf{fail}(i)&\text{otherwise,}\end{cases}\\
 \mathsf{peg}_{G}(\kw{set}(X),u,i)&=
  \begin{cases}\mathsf{ok}(i+1,\langle\kw{leaf}\langle i,i+1\rangle\rangle)
       &i<|u|\land u_i\in X,\\
    \mathsf{fail}(i)&\text{otherwise,}\end{cases}\\
 \mathsf{peg}_{G}(\kw{any},u,i)&=
  \begin{cases}\mathsf{ok}(i+1,\langle\kw{leaf}\langle i,i+1\rangle\rangle)
       &i<|u|,\\
    \mathsf{fail}(i)&i=|u|.\end{cases}
\end{aligned}
\end{equation}
\begin{equation}\label{eq:peg-sequence-choice}
\begin{aligned}
 \mathsf{peg}_{G}(r_1r_2,u,i)&=
  \begin{cases}
   \mathsf{ok}(k,\bar t_1\bar t_2)&
     \mathsf{peg}_{G}(r_1,u,i)=\mathsf{ok}(j,\bar t_1)\land
     \mathsf{peg}_{G}(r_2,u,j)=\mathsf{ok}(k,\bar t_2),\\
   \mathsf{fail}(h)&\mathsf{peg}_{G}(r_1,u,i)=\mathsf{fail}(h),\\
   \mathsf{fail}(h)&\mathsf{peg}_{G}(r_1,u,i)=\mathsf{ok}(j,\bar t_1)\land
                      \mathsf{peg}_{G}(r_2,u,j)=\mathsf{fail}(h),
  \end{cases}\\
 \mathsf{peg}_{G}(r_1/r_2,u,i)&=
  \begin{cases}
   q&\mathsf{peg}_{G}(r_1,u,i)=q=\mathsf{ok}(j,\bar t),\\
   \mathsf{peg}_{G}(r_2,u,i)&
      \mathsf{peg}_{G}(r_1,u,i)=\mathsf{fail}(h).
  \end{cases}
\end{aligned}
\end{equation}
Greedy repetition and lookahead are
\begin{equation}\label{eq:peg-repetition-lookahead}
\begin{aligned}
 \mathsf{peg}_{G}(r^{*},u,i)&=
  \begin{cases}
   \mathsf{ok}(i,\epsilon)&\mathsf{peg}_{G}(r,u,i)=\mathsf{fail}(h),\\
   \mathsf{ok}(k,\bar t\bar t')&
     \mathsf{peg}_{G}(r,u,i)=\mathsf{ok}(j,\bar t)\land j>i\land
     \mathsf{peg}_{G}(r^{*},u,j)=\mathsf{ok}(k,\bar t'),
  \end{cases}\\
 \mathsf{peg}_{G}(!r,u,i)&=
  \begin{cases}\mathsf{ok}(i,\epsilon)&
          \mathsf{peg}_{G}(r,u,i)=\mathsf{fail}(h),\\
        \mathsf{fail}(i)&\mathsf{peg}_{G}(r,u,i)=\mathsf{ok}(j,\bar t),
  \end{cases}\\
 \mathsf{peg}_{G}(\&r,u,i)&=
  \begin{cases}\mathsf{ok}(i,\epsilon)&
          \mathsf{peg}_{G}(r,u,i)=\mathsf{ok}(j,\bar t),\\
        \mathsf{fail}(i)&\mathsf{peg}_{G}(r,u,i)=\mathsf{fail}(h).
  \end{cases}
\end{aligned}
\end{equation}
Nonterminal invocation and capture are
\begin{equation}\label{eq:peg-nonterminal-capture}
\begin{aligned}
 \mathsf{peg}_{G}(N,u,i)&=\mathsf{peg}_{G}(\rho_G(N),u,i),\\
 \mathsf{peg}_{G}(\kw{capture}(n,r),u,i)&=
  \begin{cases}
   \mathsf{ok}(j,\langle\kw{node}\langle n,i,j,\bar t\rangle\rangle)&
      \mathsf{peg}_{G}(r,u,i)=\mathsf{ok}(j,\bar t),\\
   \mathsf{fail}(h)&\mathsf{peg}_{G}(r,u,i)=\mathsf{fail}(h).
  \end{cases}
\end{aligned}
\end{equation}
Undefined nonterminals are failures.  The grammar admissibility predicate is
\begin{equation}\label{eq:peg-grammar-validity}
\begin{split}
 \mathsf{pegOK}(G)\Longleftrightarrow{}&
 N_0\in\dom(\rho_G)\ \land\
 \mathsf{refsNT}(\mathsf{range}(\rho_G))\subseteq\dom(\rho_G)\ \land\\
 &\neg\exists N.\;N\mathrel{\mathsf{lead}_{G}^{+}}N\ \land\
 \forall r^{*}\preceq\mathsf{range}(\rho_G).\;\neg\mathsf{nullable}_{G}(r).
\end{split}
\end{equation}
Here $\mathsf{nullable}_{G}$ is the least predicate obtained by applying
Equations~\eqref{eq:peg-atomic-recognition}--\eqref{eq:peg-nonterminal-capture}
to empty input, and $N\mathrel{\mathsf{lead}_{G}}N'$ when $N'$ can be invoked
from $\rho_G(N)$ without consuming a character.  Thus
Equation~\eqref{eq:peg-grammar-validity} excludes left recursion and a nullable
operand of repetition; the recursive calls above consequently terminate.
The mechanization checks this relational predicate by a finite certificate.
Let $\nu_G:N\to\mathbb B$ be the computed nullable upper approximation and
let $r_G:N\to\mathbb N$ be the computed rank of leading nonterminal calls.
Write $\mathsf{nullable}^{+}_{\nu}(q)$ for the structural Boolean
approximation that interprets a nonterminal $N$ as $\nu(N)$, and write
$\mathsf{lead}^{+}_{\nu}(q)$ for the finite set of nonterminals reachable at
the head of $q$ under the same approximation.  The executable certificate is

\begin{equation}\label{eq:peg-admissibility-certificate}
\begin{split}
 \mathsf{pegCertOK}(G)\Longleftrightarrow{}&
 N_0\in\dom(\rho_G)\ \land\\
 &\forall N\in\dom(\rho_G).\ \rho_G(N)=b\Rightarrow\\
 &\quad \mathsf{refsNT}(b)\subseteq\dom(\rho_G)\ \land\\
 &\quad(\mathsf{nullable}^{+}_{\nu_G}(b)\Rightarrow\nu_G(N))\ \land\\
 &\quad\forall q^*\preceq b.\ \neg\mathsf{nullable}^{+}_{\nu_G}(q)\ \land\\
 &\quad\forall N'\in\mathsf{lead}^{+}_{\nu_G}(b).\
               r_G(N')<r_G(N).
\end{split}
\end{equation}
In the fourth line $q$ is the operand of that particular occurrence $q^*$.
Character-set
predicates are erased only while computing this structural certificate: they
cannot recognize empty input and therefore cannot affect any of its four
analyses.  The certificate checker and its semantic consequence are

\begin{equation}\label{eq:peg-certificate-soundness}
 \mathsf{pegCertOK}(G)=\kw{true}\Rightarrow\mathsf{pegOK}(G),
 \qquad
 \mathsf{pegCertOK}(G_{\NS})=\kw{true}.
\end{equation}
Thus the normalized grammar below is not merely assumed admissible; its finite
certificate is evaluated and the checker is proved sound for
Equation~\eqref{eq:peg-grammar-validity}.
Complete recognition is
\begin{equation}\label{eq:peg-complete-parse}
 \mathsf{parse}_{G}(u)=t
 \quad\Longleftrightarrow\quad
 \mathsf{pegOK}(G)\land
 \mathsf{peg}_{G}(N_0,u,0)=\mathsf{ok}(|u|,\langle t\rangle).
\end{equation}
The remainder of this subsection gives the distinguished grammar explicitly.
For compactness, the following are definitions, not additional PEG operators:
\begin{equation}\label{eq:peg-derived-notation}
 \begin{gathered}
 r?\equiv r/\epsilon,\qquad r^+\equiv rr^*,\qquad
 \lit{c_1\cdots c_n}\equiv
   \kw{char}(c_1)\cdots\kw{char}(c_n),\\
 \mathsf{triv}\equiv(\mathsf{white}/\mathsf{comment})^*,\qquad
 \mathsf{tok}(w)\equiv\mathsf{triv}\,\lit w,\qquad
 \mathsf{word}(w)\equiv\mathsf{triv}\,\lit w\,!\mathsf{idRest}.
 \end{gathered}
\end{equation}
$\mathsf{white}$ recognizes one Unicode \texttt{White\_Space} scalar.
The finite punctuation family is
\begin{equation}\label{eq:newspeak-punctuation-peg}
\begin{gathered}
 \rho_{\NS}(T_p)=\mathsf{tok}(p),\qquad p\in\mathcal P,\\
 \mathcal P=\{\texttt{:},\texttt{,},\texttt{.},\texttt{||},\texttt{=},
 \texttt{\char94},\texttt{[},\{,\texttt{(},\texttt{<},\texttt{<:},\\
 \texttt{\#},\texttt{>},\texttt{]},\},\texttt{)},\texttt{;},
 \texttt{/},\texttt{|},\texttt{::=},\texttt{<-:}\}.
\end{gathered}
\end{equation}
The longest member is selected before any of its prefixes.  This resolves the
otherwise implicit lexical priority of \texttt{||}, \texttt{::=},
\texttt{<:}, and \texttt{<-:}.

Let $L,U,D$ be the ASCII lower-case, upper-case, and decimal-digit sets,
$A=L\cup U$, and let $Q$ be the Specification's finite
\emph{specialCharacter} set.  The lexical part of $G_{\NS}$ is
\begin{equation}\label{eq:newspeak-lexical-peg}
\begin{aligned}
 \rho_{\NS}(\mathsf{idRest})&=\kw{set}(A\cup D)/\kw{char}(\texttt{\_}),\\
 \rho_{\NS}(\mathsf{rawId})&=(\kw{set}(A)/\kw{char}(\texttt{\_}))
   \mathsf{idRest}^*,\\
 \rho_{\NS}(\mathsf{id})&=\mathsf{triv}\,!\mathsf{reservedWord}
      \mathsf{rawId},\\
 \rho_{\NS}(\mathsf{rawKeyword})&=\mathsf{rawId}\lit{:},\\
 \rho_{\NS}(\mathsf{keyword})&=\mathsf{triv}\,\mathsf{rawKeyword},\\
 \rho_{\NS}(\mathsf{setterKeyword})&=
   \mathsf{triv}\,\mathsf{rawKeyword}\lit{:},\\
 \rho_{\NS}(\mathsf{keywords})&=\mathsf{rawKeyword}^+,\\
 \rho_{\NS}(\mathsf{rawBinary})&=(\kw{set}(Q)/\lit{-})\kw{set}(Q)^*,\\
 \rho_{\NS}(\mathsf{binarySelector})&=
   \mathsf{triv}\,\mathsf{rawBinary},\\
 \rho_{\NS}(\mathsf{metadataTag})&=\lit{:}\mathsf{rawId}\lit{:},\\
 \rho_{\NS}(\mathsf{commentItem})&=
   (!\lit{*)}\,\kw{any})/\mathsf{comment},\\
 \rho_{\NS}(\mathsf{comment})&=\lit{(*}\,\mathsf{metadataTag}?\\[-1mm]
   &\quad\mathsf{commentItem}^*\lit{*)}.
\end{aligned}
\end{equation}
Reserved words are recognized by the dedicated tokens
$\mathsf{selfToken}$, $\mathsf{superToken}$, $\mathsf{outerToken}$,
$\mathsf{trueToken}$, $\mathsf{falseToken}$, and $\mathsf{nilToken}$;
$\mathsf{id}$ additionally has negative lookahead for those complete words.
Thus the prose reservation is part of the grammar rather than a later
implementation convention.
The dedicated-word productions are
\begin{equation}\label{eq:newspeak-word-tokens}
\begin{aligned}
 \rho_{\NS}(\mathsf{reservedWord})&=
   \mathsf{word}(\texttt{self})/\mathsf{word}(\texttt{super})/
   \mathsf{word}(\texttt{outer})/\\[-1mm]
 &\quad\mathsf{word}(\texttt{true})/\mathsf{word}(\texttt{false})/
   \mathsf{word}(\texttt{nil}),\\
 \rho_{\NS}(\mathsf{selfToken})&=\mathsf{word}(\texttt{self}),\\
 \rho_{\NS}(\mathsf{superToken})&=\mathsf{word}(\texttt{super}),\\
 \rho_{\NS}(\mathsf{outerToken})&=\mathsf{word}(\texttt{outer}),\\
 \rho_{\NS}(\mathsf{trueToken})&=\mathsf{word}(\texttt{true}),\\
 \rho_{\NS}(\mathsf{falseToken})&=\mathsf{word}(\texttt{false}),\\
 \rho_{\NS}(\mathsf{nilToken})&=\mathsf{word}(\texttt{nil}),\\
 \rho_{\NS}(\mathsf{classToken})&=\mathsf{word}(\texttt{class}),\\
 \rho_{\NS}(\mathsf{lazyModifier})&=\mathsf{word}(\texttt{lazy}).
\end{aligned}
\end{equation}

Numeric, string, and symbol lexemes are defined by
\begin{equation}\label{eq:newspeak-literal-peg}
\begin{aligned}
 \mathsf{digits}&=\kw{set}(D)^+,\\
 \mathsf{extendedDigits}&=\kw{set}(D\cup U)^+,\\
 \mathsf{fraction}&=T_{\texttt{.}}\mathsf{digits},\\
 \mathsf{extendedFraction}&=T_{\texttt{.}}\mathsf{extendedDigits},\\
 \mathsf{exponent}&=\lit e\lit{-}?\mathsf{digits},\\
 \mathsf{decimalNum}&=\lit{-}?\mathsf{digits}
       \mathsf{fraction}?\mathsf{exponent}?,\\
 \mathsf{radixNum}&=\mathsf{digits}\lit r\lit{-}?
       \mathsf{extendedDigits}\mathsf{extendedFraction}?
       \mathsf{exponent}?,\\
 \rho_{\NS}(\mathsf{number})&=\mathsf{triv}
       (\mathsf{radixNum}/\mathsf{decimalNum}),\\
 \rho_{\NS}(\mathsf{singleString})&=
   \lit{' }((\lit{''})/(!\lit{' }\kw{any}))^*\lit{' },\\
 \rho_{\NS}(\mathsf{doubleString})&=
   \lit{"}((\lit{""})/(!\lit{"}\kw{any}))^*\lit{"},\\
 \rho_{\NS}(\mathsf{string})&=\mathsf{triv}
   (\mathsf{singleString}/\mathsf{doubleString}),\\
 \rho_{\NS}(\mathsf{symbol})&=\mathsf{triv}\,T_{\texttt{\#}}\\[-1mm]
  &\quad(\mathsf{string}/\mathsf{keywords}/
          \mathsf{rawBinary}/\mathsf{rawId}).
\end{aligned}
\end{equation}
The action in Equation~\eqref{eq:newspeak-literal-actions} below performs radix
and escape decoding and NFC normalization; recognition itself loses no source
spelling.

Optional type annotations affect parsing, tooling, and reflection but not the
untyped dynamic kernel.  Let
\begin{equation}\label{eq:newspeak-type-word-set}
\begin{split}
 \mathcal W_T=\{&\texttt{arg},\texttt{for},\texttt{generic},
 \texttt{inheritedTypeOf},\texttt{is},\texttt{message},\texttt{of},\\[-1mm]
 &\texttt{receiverType},\texttt{subtypeOf},\texttt{typeArg},\texttt{where}\}.
\end{split}
\end{equation}
Their complete PEG is
\begin{equation}\label{eq:newspeak-type-peg}
\begin{aligned}
 \rho_{\NS}(T_w)&=\mathsf{word}(w)\quad(w\in\mathcal W_T),\\
 \rho_{\NS}(\mathsf{returnType})&=T_{\texttt{\char94}}\mathsf{type},\\
 \rho_{\NS}(\mathsf{type})&=T_{\texttt{<}}\mathsf{typeExpr}T_{\texttt{>}},\\
 \rho_{\NS}(\mathsf{typePrimary})&=\mathsf{id}\mathsf{typeArguments}?,\\
 \rho_{\NS}(\mathsf{typeFactor})&=\mathsf{typePrimary}/\mathsf{blockType}/\\[-1mm]
 &\quad\mathsf{tupleType}/\mathsf{parenthesizedTypeExpression},\\
 \rho_{\NS}(\mathsf{parenthesizedTypeExpression})&=
      T_{\texttt{(}}\mathsf{typeExpr}T_{\texttt{)}},\\
 \rho_{\NS}(\mathsf{typeTerm})&=\mathsf{typeFactor}\mathsf{id}^*,\\
 \rho_{\NS}(\mathsf{typeExpr})&=\mathsf{typeTerm}\\[-1mm]
 &\quad((T_{\texttt{|}}/T_{\texttt{;}}/T_{\texttt{/}})
      \mathsf{typeExpr})?,\\
 \rho_{\NS}(\mathsf{typeArguments})&=T_{\texttt{[}}\mathsf{typeExpr}
      (T_{\texttt{,}}\mathsf{typeExpr})^*T_{\texttt{]}},\\
 \rho_{\NS}(\mathsf{tupleType})&=T_{\{}(
      \mathsf{typeExpr}(T_{\texttt{.}}\mathsf{typeExpr})^*)?T_{\}},\\
 \rho_{\NS}(\mathsf{blockArgType})&=T_{\texttt{:}}\mathsf{typeTerm},\\
 \rho_{\NS}(\mathsf{blockReturnType})&=\mathsf{typeExpr},\\
 \rho_{\NS}(\mathsf{nonEmptyBlockArgList})&=
      \mathsf{blockArgType}^+(T_{\texttt{|}}\mathsf{blockReturnType})?,\\
 \rho_{\NS}(\mathsf{blockType})&=T_{\texttt{[}}
      (\mathsf{nonEmptyBlockArgList}/\mathsf{blockReturnType}?)T_{\texttt{]}}.
\end{aligned}
\end{equation}
Type parameters and inference patterns are
\begin{equation}\label{eq:newspeak-type-pattern-peg}
\begin{aligned}
 \rho_{\NS}(\mathsf{typeParameterDecl})&=T_{\texttt{<}}T_{\texttt{[}}
      \mathsf{id}(T_{\texttt{,}}\mathsf{id})^*T_{\texttt{]}}T_{\texttt{>}},\\
 \rho_{\NS}(\mathsf{typePattern})&=T_{\texttt{<}}\mathsf{typeFormal}
      (T_{\texttt{;}}\mathsf{typeFormal})^*T_{\texttt{>}},\\
 \rho_{\NS}(\mathsf{typeFormal})&=T_{\texttt{where}}\mathsf{id}
      \mathsf{typeParamConstraint}?T_{\texttt{is}}\mathsf{inferenceClause},\\
 \rho_{\NS}(\mathsf{typeParamConstraint})&=T_{\texttt{<}}
      \mathsf{typeBoundQualifier}?\mathsf{typeExpr}T_{\texttt{>}},\\
 \rho_{\NS}(\mathsf{typeBoundQualifier})&=
      T_{\texttt{subtypeOf}}/T_{\texttt{inheritedTypeOf}},\\
 \rho_{\NS}(\mathsf{inferenceClause})&=T_{\texttt{receiverType}}/
      (\mathsf{returnType}\mathsf{returnTypeInferenceClause})/
      \mathsf{typeArgInferenceClause}/\\[-1mm]
 &\quad(T_{\texttt{arg}}\mathsf{number}
      (T_{\texttt{of}}\mathsf{msgSelector})?),\\
 \rho_{\NS}(\mathsf{returnTypeInferenceClause})&=
      T_{\texttt{of}}\mathsf{msgSelector},\\
 \rho_{\NS}(\mathsf{msgSelector})&=
      \mathsf{symbol}T_{\texttt{message}}T_{\texttt{of}}
      \mathsf{inferenceClause},\\
 \rho_{\NS}(\mathsf{typeArgInferenceClause})&=T_{\texttt{typeArg}}
      \mathsf{number}T_{\texttt{for}}T_{\texttt{generic}}\\[-1mm]
 &\quad\mathsf{symbol}T_{\texttt{of}}\mathsf{inferenceClause}.
\end{aligned}
\end{equation}

Let $\mathcal T_I$ contain every nonterminal on the left of
Equation~\eqref{eq:newspeak-type-peg} except $\mathsf{type}$, together with
every left-hand nonterminal of Equation~\eqref{eq:newspeak-type-pattern-peg}
except $\mathsf{typeParameterDecl}$ and $\mathsf{typePattern}$.  Type syntax is
retained as source metadata by
\begin{equation}\label{eq:newspeak-type-actions}
\begin{aligned}
 \mathsf{staticTypeToken}_{z}(n,w)&=
   \langle\kw{staticType},\langle n,w\rangle,z\rangle,\\
 \mathcal J_{\NS}(n,\bar v,w)&=\kw{erase}\\[-1mm]
 &\quad(n\in\mathcal T_I\cup\{T_w\mid w\in\mathcal W_T\}),\\
 \mathcal J_{\NS}(n,\bar v,w)&=
   \kw{metadata}(\mathsf{staticTypeToken}_{z}(n,w))
   \quad(n\in\{\mathsf{type},\mathsf{typeParameterDecl},
                 \mathsf{typePattern}\}).
\end{aligned}
\end{equation}

The message and precedence productions are the following PEG.  They include
the asynchronous-send token required by Specification \S5.3.7, which the
printed 0.109 grammar describes in prose but omits from its displayed
productions.
\begin{equation}\label{eq:newspeak-message-peg}
\begin{aligned}
 \rho_{\NS}(\mathsf{unarySelector})&=\mathsf{id},\\
 \rho_{\NS}(\mathsf{binaryMessage})&=
      \mathsf{binarySelector}\,\mathsf{unaryExpression},\\
 \rho_{\NS}(\mathsf{keywordMessage})&=
      (\mathsf{keyword}\,\mathsf{binaryExpression})^+,\\
 \rho_{\NS}(\mathsf{message})&=
      \mathsf{keywordMessage}/\mathsf{binaryMessage}/\mathsf{unarySelector},\\
 \rho_{\NS}(\mathsf{receiverPrimary})&=
      \mathsf{outerReceiver}/\mathsf{superToken}/\mathsf{primary},\\
 \rho_{\NS}(\mathsf{unaryExpression})&=
      \mathsf{receiverPrimary}\,\mathsf{unarySelector}^*,\\
 \rho_{\NS}(\mathsf{binaryExpression})&=
      \mathsf{unaryExpression}\,\mathsf{binaryMessage}^*,\\
 \rho_{\NS}(\mathsf{keywordExpression})&=
      \mathsf{binaryExpression}\,\mathsf{keywordMessage}?,\\
 \rho_{\NS}(\mathsf{eventualExpression})&=
      \mathsf{keywordExpression}\,T_{\texttt{<-:}}\,\mathsf{message},\\
 \rho_{\NS}(\mathsf{cascadeMessage})&=
      T_{\texttt{;}}(\mathsf{keywordMessage}/\mathsf{binaryMessage}/
                       \mathsf{unarySelector}),\\
 \rho_{\NS}(\mathsf{implicitKeywordSend})&=\mathsf{keywordMessage},\\
 \rho_{\NS}(\mathsf{nontrivialUnaryMessages})&=
      \mathsf{unarySelector}^+\mathsf{binaryMessage}^*
      \mathsf{keywordMessage}?,\\
 \rho_{\NS}(\mathsf{nontrivialBinaryMessages})&=
      \mathsf{binaryMessage}^+\mathsf{keywordMessage}?,\\
 \rho_{\NS}(\mathsf{keywordMessages})&=\mathsf{keywordMessage},\\
 \rho_{\NS}(\mathsf{nonEmptyMessages})&=
      \mathsf{nontrivialUnaryMessages}/\\[-1mm]
 &\quad\mathsf{nontrivialBinaryMessages}/\mathsf{keywordMessages},\\
 \rho_{\NS}(\mathsf{messageCascade})&=
      \mathsf{nonEmptyMessages}\,\mathsf{cascadeMessage}^*,\\
 \rho_{\NS}(\mathsf{cascadedExpression})&=
      \mathsf{primary}\,\mathsf{messageCascade}?,\\
 \rho_{\NS}(\mathsf{sendExpression})&=
      \mathsf{eventualExpression}/\mathsf{implicitKeywordSend}/
      \mathsf{cascadedExpression},\\
 \rho_{\NS}(\mathsf{expression})&=
      \mathsf{setterKeyword}?\,\mathsf{sendExpression}.
\end{aligned}
\end{equation}
The first selector of a receiver-less unary or keyword expression is retained
as an implicit send by the actions below.  Binary sends always have an
explicit left operand.  A bare \kw{outer} or \kw{super} is rejected by
Equation~\eqref{eq:surface-context-validity}; permitting it temporarily in
$\mathsf{receiverPrimary}$ keeps the precedence grammar non-left-recursive.

The literal, block, pattern, statement, and primary productions are
\begin{equation}\label{eq:newspeak-expression-peg}
\begin{aligned}
 \rho_{\NS}(\mathsf{literal})&=\mathsf{pattern}/\mathsf{number}/
   \mathsf{symbol}/\mathsf{string}/\\[-1mm]
 &\quad\mathsf{tuple}/\mathsf{trueToken}/\mathsf{falseToken}/
   \mathsf{nilToken},\\
 \rho_{\NS}(\mathsf{tuple})&=T_{\{}
   (\mathsf{expression}(T_{\texttt{.}}\mathsf{expression})^*
      T_{\texttt{.}}?)?T_{\}},\\
 \rho_{\NS}(\mathsf{blockParameter})&=T_{\texttt{:}}\mathsf{slotDecl},\\
 \rho_{\NS}(\mathsf{blockParameters})&=
      \mathsf{blockParameter}^+T_{\texttt{|}},\\
 \rho_{\NS}(\mathsf{block})&=T_{\texttt{[}}
      \mathsf{blockParameters}?\mathsf{codeBody}T_{\texttt{]}},\\
 \rho_{\NS}(\mathsf{pattern})&=T_{\texttt{<}}\mathsf{patternLiteral}T_{\texttt{>}},\\
 \rho_{\NS}(\mathsf{patternLiteral})&=\mathsf{wildcardPattern}/
      \mathsf{literalPattern}/\mathsf{keywordPattern},\\
 \rho_{\NS}(\mathsf{wildcardPattern})&=\mathsf{tok}(\texttt{\_}),\\
 \rho_{\NS}(\mathsf{literalPattern})&=\mathsf{number}/\mathsf{symbol}/
      \mathsf{string}/\mathsf{tuple},\\
 \rho_{\NS}(\mathsf{keywordPattern})&=\mathsf{keywordPatternPair}^+,\\
 \rho_{\NS}(\mathsf{keywordPatternPair})&=
      \mathsf{keyword}\,\mathsf{keywordPatternValue}?,\\
 \rho_{\NS}(\mathsf{keywordPatternValue})&=\mathsf{wildcardPattern}/
      \mathsf{literalPattern}/\\[-1mm]
 &\quad\mathsf{variablePattern}/\mathsf{nestedPattern},\\
 \rho_{\NS}(\mathsf{variablePattern})&=
      \mathsf{triv}\lit{?}\mathsf{rawId},\\
 \rho_{\NS}(\mathsf{nestedPattern})&=\mathsf{triv}\mathsf{pattern},\\
 \rho_{\NS}(\mathsf{parenthesized})&=
      T_{\texttt{(}}\mathsf{expression}T_{\texttt{)}},\\
 \rho_{\NS}(\mathsf{primary})&=\mathsf{literal}/\mathsf{block}/
      \mathsf{classDeclaration}/\\[-1mm]
 &\quad\mathsf{objectLiteral}/\mathsf{parenthesized}/\\[-1mm]
 &\quad\mathsf{selfToken}/\mathsf{unarySelector},\\
 \rho_{\NS}(\mathsf{returnStatement})&=T_{\texttt{\char94}}
      \mathsf{expression}T_{\texttt{.}}?,\\
 \rho_{\NS}(\mathsf{statements})&=\mathsf{returnStatement}/\\[-1mm]
 &\quad(\mathsf{expression}(T_{\texttt{.}}\mathsf{statements})?)/\epsilon,\\
 \rho_{\NS}(\mathsf{codeBody})&=\mathsf{slotDecls}?\mathsf{statements}.
\end{aligned}
\end{equation}

Finally, class, member, and compilation-unit syntax is
\begin{equation}\label{eq:newspeak-declaration-peg}
\begin{aligned}
 \rho_{\NS}(\mathsf{messagePattern})&=(\mathsf{keywordPatternDecl}/
      \mathsf{binaryPatternDecl}/\mathsf{unarySelector})\mathsf{returnType}?,\\
 \rho_{\NS}(\mathsf{keywordPatternDecl})&=
      (\mathsf{keyword}\mathsf{slotDecl})^+,\\
 \rho_{\NS}(\mathsf{binaryPatternDecl})&=
      \mathsf{binarySelector}\mathsf{slotDecl},\\
 \rho_{\NS}(\mathsf{accessModifier})&=
      \mathsf{tok}(\texttt{private})/\mathsf{tok}(\texttt{public})/\\[-1mm]
 &\quad\mathsf{tok}(\texttt{protected}),\\
 \rho_{\NS}(\mathsf{methodDecl})&=\mathsf{accessModifier}?
      \mathsf{messagePattern}T_{\texttt{=}}T_{\texttt{(}}
      \mathsf{codeBody}T_{\texttt{)}},\\
 \rho_{\NS}(\mathsf{slotDecl})&=\mathsf{id}\mathsf{type}?,\\
 \rho_{\NS}(\mathsf{slotDef})&=\mathsf{accessModifier}?\mathsf{slotDecl}
   (((T_{\texttt{=}}/T_{\texttt{::=}})\mathsf{expression}T_{\texttt{.}})/
      T_{\texttt{.}})?,\\
 \rho_{\NS}(\mathsf{sequentialSlots})&=T_{\texttt{|}}
      \mathsf{slotDef}^*T_{\texttt{|}},\\
 \rho_{\NS}(\mathsf{simultaneousSlots})&=T_{\texttt{||}}
      \mathsf{slotDef}^*T_{\texttt{||}},\\
 \rho_{\NS}(\mathsf{slotDecls})&=
      \mathsf{sequentialSlots}/\mathsf{simultaneousSlots},\\
 \rho_{\NS}(\mathsf{lazySlotDecl})&=
   ((\mathsf{lazyModifier}\mathsf{accessModifier}?)/\\[-1mm]
 &\quad(\mathsf{accessModifier}\mathsf{lazyModifier}))
   \mathsf{slotDecl}(T_{\texttt{=}}/T_{\texttt{::=}})\\[-1mm]
 &\quad\mathsf{expression}T_{\texttt{.}},\\
 \rho_{\NS}(\mathsf{classHeader})&=T_{\texttt{(}}\mathsf{comment}?
      \mathsf{slotDecls}?\mathsf{statements}T_{\texttt{)}},\\
 \rho_{\NS}(\mathsf{nestedClassDecl})&=
      \mathsf{accessModifier}?\mathsf{classDeclaration},\\
 \rho_{\NS}(\mathsf{instanceSide})&=T_{\texttt{(}}\\[-1mm]
 &\quad(\mathsf{nestedClassDecl}/\mathsf{lazySlotDecl}/
      \mathsf{methodDecl})^*T_{\texttt{)}},\\
 \rho_{\NS}(\mathsf{classSide})&=T_{\texttt{:}}T_{\texttt{(}}
      \mathsf{methodDecl}^*T_{\texttt{)}},\\
 \rho_{\NS}(\mathsf{classBody})&=\mathsf{classHeader}
      \mathsf{instanceSide}\mathsf{classSide}?,\\
 \rho_{\NS}(\mathsf{outerReceiver})&=\mathsf{outerToken}\mathsf{id},\\
 \rho_{\NS}(\mathsf{inheritancePrefix})&=\mathsf{outerReceiver}/
      \mathsf{selfToken}/\mathsf{superToken},\\
 \rho_{\NS}(\mathsf{inheritanceClause})&=\mathsf{inheritancePrefix}?
      \mathsf{unarySelector}\mathsf{message}?,\\
 \rho_{\NS}(\mathsf{mixinSuffix})&=
      (T_{\texttt{<:}}\mathsf{inheritanceClause})^+
      (T_{\texttt{.}}/\mathsf{classBody}),\\
 \rho_{\NS}(\mathsf{inheritanceAndBody})&=\mathsf{classBody}/\\[-1mm]
 &\quad(\mathsf{inheritanceClause}
        (\mathsf{classBody}/\mathsf{mixinSuffix})),\\
 \rho_{\NS}(\mathsf{classDeclaration})&=\mathsf{classToken}\mathsf{id}
      \mathsf{typeParameterDecl}?\\[-1mm]
 &\quad\mathsf{messagePattern}?T_{\texttt{=}}
      \mathsf{inheritanceAndBody},\\
 \rho_{\NS}(\mathsf{objectLiteral})&=
      (\mathsf{id}\mathsf{keywordMessage}?)?\mathsf{classBody},\\
 \rho_{\NS}(\mathsf{compilationUnit})&=\mathsf{id}\mathsf{string}?
      \mathsf{classDeclaration}\,\mathsf{triv}.
\end{aligned}
\end{equation}
$G_{\NS}=\langle\mathsf{compilationUnit},\rho_{\NS}\rangle$ with every
syntactic right-hand side in
Equations~\eqref{eq:newspeak-message-peg}--\eqref{eq:newspeak-declaration-peg}
wrapped in a same-named capture.  Equations~\eqref{eq:peg-derived-notation}--
\eqref{eq:newspeak-declaration-peg} are the normalized 0.109 grammar.  They
make explicit four facts left to prose or typography in the printed grammar:
Unicode whitespace, reserved-word exclusion, asynchronous sends, and the use
of class declarations, object literals, and \kw{self} as primaries.  No
unimplemented chain production is included.

\subsection{Concrete trees, metadata, and stable AST identities}
\label{sec:ast-construction}

Let $z=[i,j)$ be the span of a captured node.  Message clauses and patterns are
not expressions; their complete surface domains are
\begin{equation}\label{eq:surface-message-pattern-domain}
\begin{aligned}
 c&::=\kw{clause}_{z}(s,\bar R),\\
 p&::=\kw{wild}_{z}\mid\kw{literalPat}_{z}(R)
   \mid\kw{variablePat}_{z}(x)\mid\kw{nestedPat}_{z}(p)
   \mid\kw{keywordPat}_{z}(\langle k_i,p_i?\rangle_{i=1}^{n}).
\end{aligned}
\end{equation}
The complete executable surface AST is
\begin{equation}\label{eq:surface-expression-domain}
\begin{aligned}
 R::={}&\kw{atom}_{z}(\omega)\mid\kw{identifier}_{z}(x)
  \mid\kw{self}_{z}\mid\kw{outer}_{z}(x)\mid\kw{super}_{z}
  \mid\kw{paren}_{z}(R)\\
 &\mid\kw{implicit}_{z}(s,\bar R)
  \mid\kw{message}_{z}(R,s,\bar R)
  \mid\kw{eventual}_{z}(R,s,\bar R)
  \mid\kw{cascade}_{z}(R,\bar c_0,\bar c)\\
 &\mid\kw{setter}_{z}(x,R)\mid\kw{tuple}_{z}(\bar R)
  \mid\kw{pattern}_{z}(p)\mid\kw{closure}_{z}(\bar d,\bar g,\bar s)\\
 &\mid\kw{object}_{z}(o_h)\mid\kw{classExpr}_{z}(d)
  \mid\kw{sequence}_{z}(\bar s),\\
 s_R::={}&\kw{exprStmt}_{z}(R)\mid\kw{returnStmt}_{z}(R).
\end{aligned}
\end{equation}
Declarations and class structure are
\begin{equation}\label{eq:surface-declaration-domain}
\begin{aligned}
 d::={}&\kw{formal}_{z}(x)\mid
  \kw{slot}_{z}(\alpha?,x,m,e?)\mid
  \kw{lazySlot}_{z}(\alpha?,x,m,e)\\
 &\mid\kw{method}_{z}(\alpha?,s,\bar d,\bar g,\bar s)
  \mid\kw{nestedClass}_{z}(\alpha?,d)
  \mid\kw{classDecl}_{z}(x,s_f,\bar d,b),\\
 g::={}&\kw{sequentialGroup}_{z}(\bar d)\mid
         \kw{simultaneousGroup}_{z}(\bar d),\\
 h::={}&\kw{classStructure}_{z}(g?,\bar s,\bar d,\bar d_c),\\
 o_h::={}&\kw{implicitObjectHeader}_{z}(h)\mid
       \kw{explicitObjectHeader}_{z}(R,c,h),\\
 b::={}&\kw{defaultInheritance}_{z}(h)\mid
  \kw{explicitInheritance}_{z}(R,c,h)\mid
  \kw{mixinChain}_{z}(R,c,\langle c_i\rangle_{i=1}^{n},h?),\\
 U::={}&\kw{unit}_{z}(\lambda,\gamma?,d).
\end{aligned}
\end{equation}
Here $m\in\{\kw{immutable},\kw{mutable}\}$ records \texttt{=} versus
\texttt{::=}; $\gamma$ is the optional category string; and absent access
remains absent until elaboration supplies the protected default.  Every
declaration and expression constructor carries its captured source span and,
after Equation~\eqref{eq:stable-ast-identification}, its stable identity.

The grammar carries a finite semantic-action map
\begin{equation}\label{eq:peg-action-map}
\begin{aligned}
 \mathcal J_G:
   \mathit{CaptureName}\times\mathit{ActionValue}^{*}\times
   \mathit{SourceSlice}&\rightharpoonup\mathit{ActionResult},\\
 a&::=\kw{emit}(v)\mid\kw{erase}\mid\kw{metadata}(m).
\end{aligned}
\end{equation}
Projection returns an ordered sequence because punctuation contributes no
semantic child.  Metadata likewise leaves the semantic-child sequence but is
accumulated separately.  Put
\begin{equation}\label{eq:ast-tree-projection}
\begin{aligned}
 \mathsf{project}_{G,u}(\kw{leaf}\langle i,j\rangle)
   &=\langle\langle\kw{token}(u[i:j],[i,j))\rangle,\epsilon\rangle,\\
 \bar v&=\bar v_1\cdots\bar v_r,
 &\bar m&=\bar m_1\cdots\bar m_r,\\[-1mm]
 \mathsf{project}_{G,u}(t_k)&=\langle\bar v_k,\bar m_k\rangle
   \quad(1\le k\le r),\\
 \mathsf{project}_{G,u}(\kw{node}\langle n,i,j,\bar t\rangle)
   &=\mathsf{emitResult}
      (\mathcal J_G(n,\bar v,u[i:j]),\bar m),\\
 \mathsf{emitResult}(\kw{emit}(v),\bar m)&=\langle\langle v\rangle,\bar m\rangle,\\
 \mathsf{emitResult}(\kw{erase},\bar m)&=\langle\epsilon,\bar m\rangle,\\
 \mathsf{emitResult}(\kw{metadata}(m),\bar m)&=
      \langle\epsilon,\bar m\cdot\langle m\rangle\rangle,\\
 \mathsf{ast}_{G,u}(t)&=\langle R,\bar m\rangle
   \quad\Longleftrightarrow\quad
   \mathsf{project}_{G,u}(t)=\langle\langle R\rangle,\bar m\rangle.
\end{aligned}
\end{equation}
The metadata sequence $\bar m$ is attached by
Equation~\eqref{eq:metadata-association}; it is not an executable AST child.
Below, an equation $\mathcal J_G(\cdots)=v$ abbreviates
$\mathcal J_G(\cdots)=\kw{emit}(v)$.
To define $\mathcal J_{G_{\NS}}$ without a production-sized prose gap, first
define the message folds.  If $c_i=\kw{clause}(s_i,\bar R_i)$, then
\begin{equation}\label{eq:surface-message-folds}
\begin{aligned}
 \mathsf{asClause}_{z}(\kw{selector}(s))&=\kw{clause}_{z}(s,\epsilon),\\
 \mathsf{asClause}_{z}(c)&=c
   \quad(c=\kw{clause}_{z'}(s,\bar R)),\\
 \mathsf{flattenClauses}_{z}(\epsilon)&=\epsilon,\\
 \mathsf{flattenClauses}_{z}(v\cdot\bar v)&=
   \mathsf{asClause}_{z}(v)\cdot\mathsf{flattenClauses}_{z}(\bar v)
   \quad(v\notin\mathit{Sequence}),\\
 \mathsf{flattenClauses}_{z}(\bar c\cdot\bar v)&=
   \bar c\cdot\mathsf{flattenClauses}_{z}(\bar v)
   \quad(\bar c\in\mathit{Clause}^{*}),\\
 \mathsf{foldExplicit}_{z}(R,\epsilon)&=R,\\
 \mathsf{foldExplicit}_{z}(R,v_1\cdot\bar v)&=
   \mathsf{foldExplicit}_{z}
     (\kw{message}_{z}(R,s_1,\bar R_1),\bar v),\\[-1mm]
 &\qquad\mathsf{asClause}_{z}(v_1)=\kw{clause}(s_1,\bar R_1),\\
 \mathsf{foldImplicit}_{z}(v_1\cdot\bar v)&=
   \mathsf{foldExplicit}_{z}
     (\kw{implicit}_{z}(s_1,\bar R_1),\bar v),\\[-1mm]
 &\qquad\mathsf{asClause}_{z}(v_1)=\kw{clause}(s_1,\bar R_1),\\
 \mathsf{makeCascade}_{z}(R,\bar c_0,\epsilon)&=
   \mathsf{foldExplicit}_{z}(R,\bar c_0),\\
 \mathsf{makeCascade}_{z}(R,\bar c_0,c_1\cdot\bar c)&=
   \kw{cascade}_{z}(R,\bar c_0,c_1\cdot\bar c).
\end{aligned}
\end{equation}
The first two equations preserve unary-before-binary-before-keyword order
already imposed by Equation~\eqref{eq:newspeak-message-peg}; the third retains
one receiver and every cascade clause.

Lexical actions are
\begin{equation}\label{eq:newspeak-literal-actions}
\begin{aligned}
 \mathcal J_{\NS}(\mathsf{id},\bar v,w)&=
   \kw{name}(\mathsf{decodeId}(w)),\\
 \mathcal J_{\NS}(\mathsf{unarySelector},\langle\kw{name}(x)\rangle,w)
   &=\kw{selector}_{z}(\#x),\\
 \mathcal J_{\NS}(\mathsf{binarySelector},\bar v,w)&=
   \kw{selector}_{z}(\mathsf{decodeBinary}(w)),\\
 \mathcal J_{\NS}(\mathsf{keyword},\bar v,w)&=
   \kw{keyword}_{z}(\mathsf{decodeKeyword}(w)),\\
 \mathcal J_{\NS}(\mathsf{setterKeyword},\bar v,w)&=
   \kw{name}(\mathsf{decodeSetterName}(w)),\\
 \mathcal J_{\NS}(\mathsf{selfToken},\bar v,w)&=\kw{self}_{z},\\
 \mathcal J_{\NS}(\mathsf{superToken},\bar v,w)&=\kw{super}_{z},\\
 \mathcal J_{\NS}(\mathsf{outerReceiver},
      \langle\kw{name}(x)\rangle,w)&=\kw{outer}_{z}(x),\\
 \mathcal J_{\NS}(\mathsf{comment},\bar v,w)&=
   \kw{metadata}(\mathsf{metadataToken}_{z}(\mathsf{tag}(\bar v),
      \mathsf{payload}(w))),\\
 \mathcal J_{\NS}(\mathsf{number},\bar v,w)&=
   \kw{atom}_{z}(\mathsf{decodeNumber}(w)),\\
 \mathcal J_{\NS}(\mathsf{string},\bar v,w)&=
   \kw{atom}_{z}(\mathsf{NFC}(\mathsf{decodeString}(w))),\\
 \mathcal J_{\NS}(\mathsf{symbol},\bar v,w)&=
   \kw{atom}_{z}(\mathsf{decodeSymbol}(w)),\\
 \mathcal J_{\NS}(\mathsf{trueToken},\bar v,w)&=\kw{atom}_{z}(\kw{true}),\\
 \mathcal J_{\NS}(\mathsf{falseToken},\bar v,w)&=\kw{atom}_{z}(\kw{false}),\\
 \mathcal J_{\NS}(\mathsf{nilToken},\bar v,w)&=\kw{atom}_{z}(\kw{nil}).
\end{aligned}
\end{equation}
$\mathsf{decodeNumber}$ is the positional interpretation in Specification
\S5.1.1; it rejects a digit outside the stated radix.

Punctuation and lexical scaffolding are erased only after their enclosing
capture has used its source slice.  Pure wrapper productions pass their unique
semantic child through:
\begin{equation}\label{eq:newspeak-structural-actions}
\begin{aligned}
 \mathcal E_{\NS}={}&\{T_p\mid p\in\mathcal P\}\cup\{\mathsf{white},
   \mathsf{reservedWord},\mathsf{idRest},\mathsf{rawId},\\[-1mm]
 &\mathsf{rawKeyword},\mathsf{keywords},\mathsf{rawBinary},
   \mathsf{metadataTag},\mathsf{commentItem},\\[-1mm]
 &\mathsf{classToken},\mathsf{lazyModifier},\mathsf{outerToken}\},\\
 \mathcal I_{\NS}={}&\{\mathsf{literal},\mathsf{receiverPrimary},
   \mathsf{sendExpression},\\[-1mm]
 &\mathsf{slotDecls},\mathsf{blockParameter},\mathsf{inheritancePrefix},
   \mathsf{patternLiteral},\\[-1mm]
 &\mathsf{keywordPatternValue}\},\\
 \mathcal J_{\NS}(n,\bar v,w)&=\kw{erase}
      \quad(n\in\mathcal E_{\NS}),\\
 \mathcal J_{\NS}(n,\langle v\rangle,w)&=v
      \quad(n\in\mathcal I_{\NS}).
\end{aligned}
\end{equation}

Sequence normalization used by the expression actions is
\begin{equation}\label{eq:newspeak-action-sequence-helpers}
\begin{aligned}
 \mathsf{orEmpty}(\kw{none})&=\epsilon,\\
 \mathsf{orEmpty}(v)&=\langle v\rangle\quad(v\ne\kw{none}),\\
 \mathsf{orEmptySeq}(\kw{none})&=\epsilon,\\
 \mathsf{orEmptySeq}(\bar v)&=\bar v\quad(\bar v\ne\kw{none}),\\
 \mathsf{asStatement}(R)&=\kw{exprStmt}(R)
      \quad(R\in\mathit{SurfaceExpression}),\\
 \mathsf{asStatement}(s_R)&=s_R
      \quad(s_R\in\mathit{SurfaceStatement}),\\
 \mathsf{flattenStatements}(\epsilon)&=\epsilon,\\
 \mathsf{flattenStatements}(v\cdot\bar v)&=
      \mathsf{asStatement}(v)\cdot\mathsf{flattenStatements}(\bar v)
      \quad(v\notin\mathit{Sequence}),\\
 \mathsf{flattenStatements}(\kw{sequence}(\bar s)\cdot\bar v)&=
      \bar s\cdot\mathsf{flattenStatements}(\bar v),\\
\end{aligned}
\end{equation}
The message and expression actions are exhaustive:
\begin{equation}\label{eq:newspeak-expression-actions}
\begin{aligned}
 \mathcal J_{\NS}(\mathsf{binaryMessage},
      \langle\kw{selector}(s),R\rangle,w)
   &=\kw{clause}_{z}(s,\langle R\rangle),\\
 \mathcal J_{\NS}(\mathsf{keywordMessage},
      \langle\kw{keyword}(k_i),R_i\rangle_{i=1}^{n},w)
   &=\kw{clause}_{z}(k_1\cdots k_n,\langle R_i\rangle_{i=1}^{n}),\\
 \mathcal J_{\NS}(\mathsf{message},
      \langle\kw{selector}(s)\rangle,w)&=\kw{clause}_{z}(s,\epsilon),\\
 \mathcal J_{\NS}(\mathsf{message},\langle c\rangle,w)&=c
      \quad(c=\kw{clause}_{z'}(s,\bar R)),\\
 \mathcal J_{\NS}(\mathsf{cascadeMessage},
      \langle\kw{selector}(s)\rangle,w)&=\kw{clause}_{z}(s,\epsilon),\\
 \mathcal J_{\NS}(\mathsf{cascadeMessage},\langle c\rangle,w)&=c,\\
\end{aligned}
\end{equation}
The productions that normalize to an ordered sequence of message clauses are
\begin{equation}\label{eq:newspeak-message-sequence-names}
 \mathcal N_M=\{\mathsf{nontrivialUnaryMessages},
   \mathsf{nontrivialBinaryMessages},
   \mathsf{keywordMessages},\mathsf{nonEmptyMessages}\}.
\end{equation}
\begin{equation}\label{eq:newspeak-send-fold-actions}
\begin{aligned}
 \mathcal J_{\NS}(\mathsf{primary},
      \langle\kw{selector}_{z'}(\#x)\rangle,w)&=\kw{identifier}_{z}(x),\\
 \mathcal J_{\NS}(\mathsf{primary},
      \langle\kw{classDecl}_{z'}(\bar v)\rangle,w)&=
      \kw{classExpr}_{z}(\kw{classDecl}_{z'}(\bar v)),\\
 \mathcal J_{\NS}(\mathsf{primary},\langle R\rangle,w)&=R
      \quad\text{for every other primary},\\
 \mathcal J_{\NS}(\mathsf{unaryExpression},\langle R,\bar c\rangle,w)
   &=\mathsf{foldExplicit}_{z}(R,\bar c),\\
 \mathcal J_{\NS}(\mathsf{binaryExpression},\langle R,\bar c\rangle,w)
   &=\mathsf{foldExplicit}_{z}(R,\bar c),\\
 \mathcal J_{\NS}(\mathsf{keywordExpression},\langle R,c?\rangle,w)
   &=\mathsf{foldExplicit}_{z}(R,c?),\\
 \mathcal J_{\NS}(\mathsf{eventualExpression},\langle R,c\rangle,w)
   &=\kw{eventual}_{z}(R,\selector(c),\mathsf{arguments}(c)),\\
 \mathcal J_{\NS}(\mathsf{implicitKeywordSend},\langle c\rangle,w)
   &=\kw{implicit}_{z}(\selector(c),\mathsf{arguments}(c)),\\
 \mathcal J_{\NS}(n,\bar v,w)&=
      \mathsf{flattenClauses}_{z}(\bar v)
      \quad(n\in\mathcal N_M),\\
 \mathcal J_{\NS}(\mathsf{messageCascade},
      \langle\bar c_0,\bar c\rangle,w)
   &=\kw{cascadeParts}_{z}(\bar c_0,\bar c),\\
 \mathcal J_{\NS}(\mathsf{cascadedExpression},\langle R\rangle,w)&=R,\\
 \mathcal J_{\NS}(\mathsf{cascadedExpression},
      \langle R,\kw{cascadeParts}(\bar c_0,\bar c)\rangle,w)
   &=\mathsf{makeCascade}_{z}(R,\bar c_0,\bar c),\\
\end{aligned}
\end{equation}
\begin{equation}\label{eq:newspeak-compound-expression-actions}
\begin{aligned}
 \mathcal J_{\NS}(\mathsf{expression},
      \langle\kw{name}(x),R\rangle,w)&=\kw{setter}_{z}(x,R),\\
 \mathcal J_{\NS}(\mathsf{expression},\langle R\rangle,w)&=R,\\
 \mathcal J_{\NS}(\mathsf{parenthesized},\langle R\rangle,w)&=\kw{paren}_{z}(R),\\
 \mathcal J_{\NS}(\mathsf{tuple},\bar R,w)&=\kw{tuple}_{z}(\bar R),\\
 \mathcal J_{\NS}(\mathsf{blockParameters},\bar d,w)&=\bar d,\\
 \mathcal J_{\NS}(\mathsf{codeBody},
      \langle g?,\kw{sequence}(\bar s)\rangle,w)
   &=\kw{body}_{z}(g?,\bar s),\\
 \mathcal J_{\NS}(\mathsf{block},
      \langle\bar d?,\kw{body}(g?,\bar s)\rangle,w)
   &=\kw{closure}_{z}(\mathsf{orEmptySeq}(\bar d?),
       \mathsf{orEmpty}(g?),\bar s),\\
 \mathcal J_{\NS}(\mathsf{returnStatement},\langle R\rangle,w)
   &=\kw{returnStmt}_{z}(R),\\
 \mathcal J_{\NS}(\mathsf{statements},\bar v,w)&=
      \kw{sequence}_{z}(\mathsf{flattenStatements}(\bar v)).
\end{aligned}
\end{equation}
A unary selector in primary position becomes $\kw{identifier}_{z}(x)$;
Equation~\eqref{eq:expression-elaboration} later makes it an implicit nullary
send.  An implicit keyword message instead takes the explicit
$\mathsf{implicitKeywordSend}$ branch.  Thus neither form introduces a
variable-read node.

Pattern actions are
\begin{equation}\label{eq:newspeak-pattern-actions}
\begin{aligned}
 \mathcal J_{\NS}(\mathsf{wildcardPattern},\epsilon,w)&=\kw{wild}_{z},\\
 \mathcal J_{\NS}(\mathsf{literalPattern},\langle R\rangle,w)&=
      \kw{literalPat}_{z}(R),\\
 \mathcal J_{\NS}(\mathsf{variablePattern},\langle\kw{name}(x)\rangle,w)&=
      \kw{variablePat}_{z}(x),\\
 \mathcal J_{\NS}(\mathsf{nestedPattern},\langle p\rangle,w)&=
      \kw{nestedPat}_{z}(p),\\
 \mathcal J_{\NS}(\mathsf{keywordPatternPair},
      \langle\kw{keyword}(k),p?\rangle,w)&=\langle k,p?\rangle,\\
 \mathcal J_{\NS}(\mathsf{keywordPattern},\bar v,w)&=
      \kw{keywordPat}_{z}(\bar v),
      \quad\bar v=\langle k_i,p_i?\rangle_{i=1}^{n},\\
 \mathcal J_{\NS}(\mathsf{pattern},\langle p\rangle,w)&=\kw{pattern}_{z}(p).
\end{aligned}
\end{equation}

For an inheritance clause define
\begin{equation}\label{eq:newspeak-inheritance-actions}
\begin{aligned}
 \mathsf{inheritanceReceiver}_{z}(\kw{none},\kw{selector}(s))
   &=\kw{implicit}_{z}(s,\epsilon),\\
 \mathsf{inheritanceReceiver}_{z}(R,\kw{selector}(s))
   &=\kw{message}_{z}(R,s,\epsilon),\\
 \mathsf{initializerClause}_{z}(\kw{none})&=
   \kw{clause}_{z}(\#\kw{new},\epsilon),\\
 \mathsf{initializerClause}_{z}(c)&=c,\\
 R_I&=\mathsf{inheritanceReceiver}_{z}(R?,\kw{selector}(s)),\\
 c_I&=\mathsf{initializerClause}_{z}(c?),\\
 \mathcal J_{\NS}(\mathsf{inheritanceClause},
      \langle R?,\kw{selector}(s),c?\rangle,w)
   &=\kw{inheritanceHead}_{z}(R_I,c_I).
\end{aligned}
\end{equation}
For compactness in the remaining action equations, define the non-overloaded
abbreviation
\begin{equation}\label{eq:newspeak-action-abbreviation}
 \mathsf{act}_{[i,j)}(n,\bar v,u[i:j])
   =\mathcal J_{\NS}(n,\bar v,u[i:j]).
\end{equation}
The complete declaration actions are
\begin{equation}\label{eq:newspeak-declaration-actions}
\begin{aligned}
 N_I&=\mathsf{inheritanceAndBody},\\
 \mathsf{sel}(\kw{patDecl}(s,\bar d))&=s,\\
 \mathsf{formals}(\kw{patDecl}(s,\bar d))&=\bar d,\\
 \mathsf{factorySel}(\kw{none})&=\#\kw{new},\\
 \mathsf{factorySel}(\kw{patDecl}(s,\bar d))&=s,\\
 \mathsf{factoryFormals}(\kw{none})&=\epsilon,\\
 \mathsf{factoryFormals}(\kw{patDecl}(s,\bar d))&=\bar d,\\
 \mathsf{act}_{z}(\mathsf{keywordPatternDecl},
      \langle\kw{keyword}(k_i),d_i\rangle_{i=1}^{n},w)&=\\[-1mm]
 &\quad\kw{patDecl}_{z}(k_1\cdots k_n,\langle d_i\rangle_{i=1}^{n}),\\
 \mathsf{act}_{z}(\mathsf{binaryPatternDecl},
      \langle\kw{selector}(s),d\rangle,w)&=
      \kw{patDecl}_{z}(s,\langle d\rangle),\\
 \mathsf{act}_{z}(\mathsf{messagePattern},
      \langle\kw{selector}(s)\rangle,w)&=\kw{patDecl}_{z}(s,\epsilon),\\
 \mathsf{act}_{z}(\mathsf{messagePattern},\langle p\rangle,w)&=p
      \quad(p=\kw{patDecl}_{z'}(s,\bar d)),\\
 \mathsf{act}_{z}(\mathsf{accessModifier},\bar v,w)&=
      \mathsf{decodeAccess}(w),\\
 \mathsf{act}_{z}(\mathsf{slotDecl},\langle\kw{name}(x)\rangle,w)&=
      \kw{formal}_{z}(x),\\
 \mathsf{act}_{z}(\mathsf{slotDef},\langle\alpha?,d,R?\rangle,w)&=\\[-1mm]
 &\quad\kw{slot}_{z}(\alpha?,\mathsf{name}(d),
      \mathsf{mutability}(w),R?),\\
 \mathsf{act}_{z}(\mathsf{sequentialSlots},\bar d,w)&=
      \kw{sequentialGroup}_{z}(\bar d),\\
 \mathsf{act}_{z}(\mathsf{simultaneousSlots},\bar d,w)&=
      \kw{simultaneousGroup}_{z}(\bar d),\\
 \mathsf{act}_{z}(\mathsf{lazySlotDecl},
      \langle\alpha?,d,R\rangle,w)&=\\[-1mm]
 &\quad\kw{lazySlot}_{z}(\alpha?,\mathsf{name}(d),
      \mathsf{mutability}(w),R),\\
 \mathsf{act}_{z}(\mathsf{methodDecl},
      \langle\alpha?,p,\kw{body}(g?,\bar s)\rangle,w)&=\\[-1mm]
 &\quad\kw{method}_{z}(\alpha?,\mathsf{sel}(p),
      \mathsf{formals}(p),\mathsf{orEmpty}(g?),\bar s),\\
 \mathsf{act}_{z}(\mathsf{nestedClassDecl},\langle\alpha?,d\rangle,w)&=
      \kw{nestedClass}_{z}(\alpha?,d),\\
 \mathsf{act}_{z}(\mathsf{classHeader},
      \langle g?,\kw{sequence}(\bar s)\rangle,w)&=
      \kw{header}_{z}(g?,\bar s),\\
 \mathsf{act}_{z}(\mathsf{instanceSide},\bar d,w)&=\bar d,\\
 \mathsf{act}_{z}(\mathsf{classSide},\bar d_c,w)&=\bar d_c,\\
 \mathsf{act}_{z}(\mathsf{classBody},
      \langle\kw{header}(g?,\bar s),\bar d,\bar d_c?\rangle,w)&=\\[-1mm]
 &\quad\kw{classStructure}_{z}
      (g?,\bar s,\bar d,\mathsf{orEmptySeq}(\bar d_c?)),\\
 \mathsf{act}_{z}(\mathsf{mixinSuffix},\langle\bar h,h?\rangle,w)&=
      \kw{mixinTail}_{z}(\bar h,h?),\\
 \mathsf{act}_{z}(N_I,\langle h\rangle,w)&=
      \kw{defaultInheritance}_{z}(h),\\
 \mathsf{act}_{z}(N_I,
      \langle\kw{inheritanceHead}(R,c),h\rangle,w)&=
      \kw{explicitInheritance}_{z}(R,c,h),\\
 h_I&=\kw{inheritanceHead}(R,c),\\
 t_I&=\kw{mixinTail}(\bar c,h?),\\
 \mathsf{act}_{z}(N_I,\langle h_I,t_I\rangle,w)&=\\[-1mm]
 &\quad\kw{mixinChain}_{z}(R,c,\bar c,h?),\\
 \mathsf{act}_{z}(\mathsf{classDeclaration},
      \langle\kw{name}(x),p?,b\rangle,w)&=
      \kw{classDecl}_{z}(x,\mathsf{factorySel}(p?),
         \mathsf{factoryFormals}(p?),b),\\
 \mathsf{act}_{z}(\mathsf{objectLiteral},\langle h\rangle,w)&=
      \kw{object}_{z}(\kw{implicitObjectHeader}_{z}(h)),\\
 R_o&=\kw{identifier}_{z}(x),\\
 c_o&=\mathsf{initializerClause}_{z}(c?),\\
 o_h&=\kw{explicitObjectHeader}_{z}(R_o,c_o,h),\\
 \mathsf{act}_{z}(\mathsf{objectLiteral},
      \langle\kw{name}(x),c?,h\rangle,w)&=\\[-1mm]
 &\quad\kw{object}_{z}(o_h),\\
 \mathsf{act}_{z}(\mathsf{compilationUnit},
      \langle\lambda,\gamma?,d\rangle,w)&=\\[-1mm]
 &\quad\kw{unit}_{z}(\lambda,\gamma?,d).
\end{aligned}
\end{equation}
$\mathsf{mutability}(w)$ is immutable exactly when the slot definition's
source slice $w$ contains \texttt{=} rather than \texttt{::=}; either
assignment operator makes a lazy slot mutable or immutable in the same way,
and omission is mutable.
The inheritance actions construct exactly the three $b$ alternatives in
Equation~\eqref{eq:surface-declaration-domain}, in source order.  All
punctuation, trivia, and wrapper productions have either
the discard action or the unique-child action.  Equations
\eqref{eq:newspeak-literal-actions}--\eqref{eq:newspeak-declaration-actions},
together with those two structural actions, define $\mathcal J_{G_{\NS}}$ on
every successful captured production; there is no residual action case.

A metadata token is $m=\langle\tau,w,[i,j)\rangle$, where $\tau$ is its tag
and $w$ its uninterpreted payload; an ordinary comment has tag $\kw{comment}$.
Let $\mathsf{triviaOnly}_{u}(i,j)$ mean that $u[i:j]$ contains only whitespace
and comments, and define
\begin{equation}\label{eq:metadata-trivia-gap}
 \mathsf{triviaBetween}_{u}(X,Y)\Longleftrightarrow
 \mathsf{end}(X)\le\mathsf{start}(Y)\land
 \mathsf{triviaOnly}_{u}(\mathsf{end}(X),\mathsf{start}(Y)).
\end{equation}
The association relation is
\begin{equation}\label{eq:metadata-association}
\begin{aligned}
 m\mathrel{\mathsf{before}_{u}}R
 &\Longleftrightarrow
   \mathsf{triviaBetween}_{u}(m,R)\land\\[-1mm]
 &\qquad R\in\mathit{Expression}\uplus\mathit{FormalDeclaration}
                    \uplus\mathit{SlotDeclaration},\\
 m\mathrel{\mathsf{afterPattern}_{u}}R
 &\Longleftrightarrow
   \mathsf{triviaBetween}_{u}(\mathsf{pattern}(R),m)\land\\[-1mm]
 &\qquad R\in\mathit{MethodDeclaration}
                   \uplus\mathit{FactoryDeclaration},\\
 m\mathrel{\mathsf{afterHeader}_{u}}R
 &\Longleftrightarrow
   \mathsf{triviaBetween}_{u}(\mathsf{header}(R),m)\land
   R\in\mathit{ClassDeclaration},\\
 m\mathrel{\mathsf{staticTypeOf}_{u}}R
 &\Longleftrightarrow
   \mathsf{tag}(m)=\kw{staticType}\land
   \mathsf{span}(m)\subseteq\mathsf{span}(R)\land\\[-1mm]
 &\qquad R=\min_{\subseteq}\{R'\mid
      \mathsf{span}(m)\subseteq\mathsf{span}(R')\land
      R'\in\mathit{TypedDeclaration}\}.
\end{aligned}
\end{equation}
The applicable relation is selected by node kind in displayed order.  Metadata
attachment and its insertion into the tree are the functions
\begin{equation}\label{eq:metadata-attachment}
\begin{aligned}
 \mathsf{related}_{u}(m,R)&=
  \begin{cases}
   m\mathrel{\mathsf{staticTypeOf}_{u}}R&\mathsf{tag}(m)=\kw{staticType},\\
   m\mathrel{\mathsf{afterHeader}_{u}}R
      &R\in\mathit{ClassDeclaration},\\
   m\mathrel{\mathsf{afterPattern}_{u}}R
      &R\in\mathit{MethodDeclaration}\uplus\mathit{FactoryDeclaration},\\
   m\mathrel{\mathsf{before}_{u}}R&\text{otherwise,}
  \end{cases}\\
 \mathsf{metadata}_{u,\bar m}(R)&=
   \mathsf{sourceOrder}(\langle m\in\bar m\mid
      \mathsf{related}_{u}(m,R)\rangle),\\
 \mathsf{attachMetadata}_{u}(R,\bar m)&=
   R[\mathsf{meta}(n)\mapsto\mathsf{metadata}_{u,\bar m}(n)
       \mid n\in\mathsf{nodes}(R)].
\end{aligned}
\end{equation}
Punctuation, whitespace, and metadata tokens therefore do not become
executable AST children.

Parsed nodes next receive identities.  Let
$\kappa:\mathit{NodeKey}\rightharpoonup
 (\mathit{DeclId}\uplus\mathit{SiteId})$ be an optional identity-retention map,
where a key contains a source-unit identity, structural path, constructor kind,
and declared name when one exists.  Define
\begin{equation}\label{eq:stable-ast-identification}
 \mathsf{identify}_{\kappa}(R)=\langle S,\kappa'\rangle
 \quad\Longleftrightarrow\quad
 \begin{cases}
  \mathsf{id}_{S}(n)=\kappa(\mathsf{key}_{R}(n))
    &\mathsf{key}_{R}(n)\in\dom(\kappa),\\
  \mathsf{id}_{S}(n)\text{ is fresh of the required sort}
    &\mathsf{key}_{R}(n)\notin\dom(\kappa),
 \end{cases}
\end{equation}
for every declaration or expression node $n$, with $\kappa'$ extending
$\kappa$ by the chosen keys.  The result must satisfy
\begin{equation}\label{eq:stable-ast-identity-validity}
 \mathsf{idsOK}(S)\Longleftrightarrow
 \mathsf{id}_{S}\text{ is injective}\ \land\
 \mathsf{id}_{S}(n)\in
 \begin{cases}\mathit{DeclId}&n\in\mathit{DeclarationNode}(S),\\
               \mathit{SiteId}&n\in\mathit{ExpressionNode}(S).
 \end{cases}
\end{equation}
An initial compilation uses $\kappa=\varnothing$.  A source-preserving editor
may supply a retained map; mirror changes that directly edit identified AST
nodes preserve identities according to Equation~\eqref{eq:source-patch} and do
not reparse text.

\subsection{Binding, elaboration, and static well-formedness}
\label{sec:static-elaboration}

An elaboration context is
\begin{equation}\label{eq:elaboration-context}
 \Gamma=\langle K,q,D_0,\bar B,r_t,h_t\rangle,
\end{equation}
where $K$ is the enclosing class-body declaration or $\kw{top}$, $q$ is the
enclosing activation declaration or $\kw{none}$, $D_0$ is the immediately
enclosing named class or $\kw{none}$, $\bar B$ is the innermost-first sequence
of lexical binder maps, and $r_t,h_t$ record whether local and non-local return
are permitted.  Lexical binding is the partial function
\begin{equation}\label{eq:static-binding-function}
 \mathsf{bind}_{\Gamma}(x)=d
 \quad\Longleftrightarrow\quad
 j=\min\{i\mid x\in\dom(B_i)\}\land d=B_j(x),
 \qquad \bar B=\langle B_0,\ldots,B_m\rangle.
\end{equation}
An absent set on the right makes the function undefined.  Class-name binding
uses the same definition restricted to class declarations; write
$\mathsf{classBind}_{\Gamma}(x)=D_t$.
The total annotation used by an ordinary implicit send is
\begin{equation}\label{eq:implicit-binding-annotation}
 \mathsf{bindingAnnotation}_{\Gamma}(x)=
 \begin{cases}
  \mathsf{bind}_{\Gamma}(x)&\mathsf{bind}_{\Gamma}(x)\ne\undef,\\
  \kw{none}&\mathsf{bind}_{\Gamma}(x)=\undef.
 \end{cases}
\end{equation}

Parentheses do not erase receiver authority.  Define
\begin{equation}\label{eq:surface-receiver-kind}
\begin{aligned}
 \mathsf{receiverKind}_{\Gamma}((R))&=\mathsf{receiverKind}_{\Gamma}(R),\\
 \mathsf{receiverKind}_{\Gamma}(\kw{self})&=\kw{selfReceiver}(D_0),\\
 \mathsf{receiverKind}_{\Gamma}(\kw{outer}(x))
   &=\kw{outerReceiver}(\mathsf{classBind}_{\Gamma}(x),D_0),\\
 \mathsf{receiverKind}_{\Gamma}(\kw{super})&=\kw{superReceiver},\\
 \mathsf{receiverKind}_{\Gamma}(R)&=\kw{ordinaryReceiver}
   \quad\text{otherwise.}
\end{aligned}
\end{equation}
For already elaborated arguments $\bar e$, put
$r_k=\mathsf{receiverKind}_{\Gamma}(R)$.  Send construction is the distinct
partial function
\begin{equation}\label{eq:elaborated-send-construction}
\begin{aligned}
 \mathsf{makeSend}_{\Gamma}(R,s,\bar e)&=
 \begin{cases}
  \kw{selfSend}(s,\bar e,D_0)&
    r_k=\kw{selfReceiver}(D_0),\\
  \kw{outerSend}(s,\bar e,D_t,D_0)&
    r_k=\kw{outerReceiver}(D_t,D_0),\\
  \kw{superSend}(s,\bar e)&
    r_k=\kw{superReceiver},\\
  \kw{send}(e_R,s,\bar e)&
    r_k=\kw{ordinaryReceiver}\land
      \mathsf{elabExpr}_{\Gamma}(R)=e_R.
 \end{cases}
\end{aligned}
\end{equation}
It is undefined if $D_0=\kw{none}$ in a self or outer case, if the named outer
class does not bind, or if super occurs outside a method activation.  Notice
again that the super case receives no declaration annotation.

Use distinct functions for a single expression, an expression list, a
statement sequence, a declaration, and a complete unit:
\begin{equation}\label{eq:elaboration-function-signatures}
\begin{aligned}
 \mathsf{elabExpr}_{\Gamma}&:\mathit{SurfaceExpression}\rightharpoonup
     \mathit{CoreExpression},\\
 \mathsf{elabExprList}_{\Gamma}&:\mathit{SurfaceExpression}^{*}
     \rightharpoonup\mathit{CoreExpression}^{*},\\
 \mathsf{elabStatements}_{\Gamma}&:\mathit{SurfaceStatement}^{*}
     \rightharpoonup\mathit{CoreStatement}^{*},\\
 \mathsf{elabDecl}_{\Gamma}&:\mathit{SurfaceDeclaration}
     \rightharpoonup\mathit{DerivedDeclaration},\\
 \mathsf{deriveProgram}&:\mathit{IdentifiedUnit}\rightharpoonup
     (\mathit{Program}\times\mathit{CoreExpression}).
\end{aligned}
\end{equation}
The list and optional functions, selector annotations, and activation contexts
are defined by
\begin{equation}\label{eq:elaboration-helper-functions}
\begin{aligned}
 \mathsf{elabExprList}_{\Gamma}(\epsilon)&=\epsilon,\\
 \mathsf{elabExprList}_{\Gamma}(R\cdot\bar R)&=
   \mathsf{elabExpr}_{\Gamma}(R)\cdot
   \mathsf{elabExprList}_{\Gamma}(\bar R),\\
 \mathsf{elabOptionalExpr}_{\Gamma}(\kw{none})&=\kw{none},\\
 \mathsf{elabOptionalExpr}_{\Gamma}(R)&=\mathsf{elabExpr}_{\Gamma}(R),\\
 \mathsf{elabOptionalPattern}_{\Gamma}(\kw{none})&=\_,\\
 \mathsf{elabOptionalPattern}_{\Gamma}(p)&=\mathsf{elabPattern}_{\Gamma}(p),\\
 \mathsf{selectorBindingAnnotation}_{\Gamma}(s)&=
  \begin{cases}
   \mathsf{bindingAnnotation}_{\Gamma}(x)&s=\#x,\\
   \kw{none}&\mathsf{arity}(s)>0.
  \end{cases}
\end{aligned}
\end{equation}
For a declaration sequence, let
\begin{equation}\label{eq:elaboration-context-extension}
\begin{aligned}
 \mathsf{groupDecls}(\epsilon)&=\epsilon,\\
 \mathsf{groupDecls}(g\cdot\bar g)&=
     \mathsf{members}(g)\cdot\mathsf{groupDecls}(\bar g),\\
 \mathsf{binder}(\bar d)&=
     \{\mathsf{name}(d)\mapsto\mathsf{id}(d)\mid d\in\bar d\},\\
 \mathsf{returnFlags}_{\Gamma}(q)&=
  \begin{cases}
   \langle\mathsf{true},\mathsf{false}\rangle&q\in\mathit{MethodDeclId},\\
   \langle\mathsf{false},\Gamma.r_t\lor\Gamma.h_t\rangle
       &q\in\mathit{ClosureDeclId},
  \end{cases}\\
 \mathsf{activationContext}_{\Gamma}(q,\bar d,\bar g)&=
   \langle\Gamma.K,q,\Gamma.D_0,
     \mathsf{binder}(\bar d\cdot\mathsf{groupDecls}(\bar g))
       \cdot\Gamma.\bar B,r_q,h_q\rangle,\\[-1mm]
 &\qquad\langle r_q,h_q\rangle=\mathsf{returnFlags}_{\Gamma}(q).
\end{aligned}
\end{equation}
Thus $\Gamma_q=\mathsf{activationContext}_{\Gamma}(q,\bar d,\bar g)$ in every
closure or method case below.  The atomic and send cases are
\begin{equation}\label{eq:expression-elaboration}
\begin{aligned}
 \mathsf{elabExpr}_{\Gamma}(\kw{identifier}_{z}(x))
   &=\kw{implicit}(\#x,\epsilon,
      \mathsf{bindingAnnotation}_{\Gamma}(x)),\\
 \mathsf{elabExpr}_{\Gamma}(\kw{implicit}_{z}(s,\bar R))
   &=\kw{implicit}(s,\mathsf{elabExprList}_{\Gamma}(\bar R),
      \mathsf{selectorBindingAnnotation}_{\Gamma}(s)),\\
 \mathsf{elabExpr}_{\Gamma}(\kw{message}_{z}(R,s,\bar R'))
   &=\mathsf{makeSend}_{\Gamma}
       (R,s,\mathsf{elabExprList}_{\Gamma}(\bar R')),\\
 \mathsf{elabExpr}_{\Gamma}(\kw{eventual}_{z}(R,s,\bar R'))
   &=\kw{asyncSend}(\mathsf{elabExpr}_{\Gamma}(R),s,
                    \mathsf{elabExprList}_{\Gamma}(\bar R')),\\
 \mathsf{elabExpr}_{\Gamma}(\kw{atom}_{z}(\omega))&=\kw{atom}(\omega),\\
 \mathsf{elabExpr}_{\Gamma}(\kw{self}_{z})&=\kw{self},\\
 \mathsf{elabExpr}_{\Gamma}(\kw{paren}_{z}(R))&=\mathsf{elabExpr}_{\Gamma}(R),\\
 \mathsf{elabExpr}_{\Gamma}(\kw{tuple}_{z}(\bar R))
   &=\kw{tuple}(\mathsf{elabExprList}_{\Gamma}(\bar R)),\\
 \mathsf{elabExpr}_{\Gamma}(\kw{pattern}_{z}(p))
   &=\kw{pattern}(\mathsf{elabPattern}_{\Gamma}(p),
       \mathsf{bindingAnnotation}_{\Gamma}(\kw{Pattern})),\\
 \mathsf{elabExpr}_{\Gamma}(\kw{setter}_{z}(x,R))
   &=\mathsf{elabExpr}_{\Gamma}(\mathsf{setterExpansion}_{z}(x,R)),\\
 \mathsf{elabExpr}_{\Gamma}(\kw{cascade}_{z}(R,\bar c_0,\bar c))
   &=\mathsf{elabExpr}_{\Gamma}
      (\mathsf{cascadeExpansion}_{z}(R,\bar c_0,\bar c)).
\end{aligned}
\end{equation}
Here $\mathsf{selectorBindingAnnotation}_{\Gamma}$ is
Equation~\eqref{eq:implicit-binding-annotation} for a unary activation member
and is $\kw{none}$ otherwise.  Tuple and pattern nodes are retained because
their dynamic expansions are Equations~\eqref{eq:tuple-expansion} and
\eqref{eq:pattern-components}--\eqref{eq:variable-pattern-parameter}.

Pattern elaboration recursively elaborates only embedded expressions:
\begin{equation}\label{eq:pattern-elaboration}
\begin{aligned}
 \mathsf{elabPattern}_{\Gamma}(\kw{wild}_{z})&=\_,\\
 \mathsf{elabPattern}_{\Gamma}(\kw{literalPat}_{z}(R))&=
      \mathsf{elabExpr}_{\Gamma}(R),\\
 \mathsf{elabPattern}_{\Gamma}(\kw{variablePat}_{z}(x))&=?x_z,\\
 \mathsf{elabPattern}_{\Gamma}(\kw{nestedPat}_{z}(p))&=
      \langle\mathsf{elabPattern}_{\Gamma}(p)\rangle,\\
 \mathsf{elabPattern}_{\Gamma}
    (\kw{keywordPat}_{z}(\langle k_i,p_i?\rangle_i))&=
      \langle k_i,
        \mathsf{elabOptionalPattern}_{\Gamma}(p_i?)\rangle_i.
\end{aligned}
\end{equation}
An absent keyword component denotes a wildcard.  The resulting forms are
exactly the arguments of Equations~\eqref{eq:pattern-components}--
\eqref{eq:keyword-pattern-expansion}.

The setter expansion, with fresh $p\ne x$, is
\begin{equation}\label{eq:setter-expansion}
\begin{aligned}
 \mathsf{setterBody}_{z}(x,p)&=
   \langle\kw{exprStmt}(\kw{implicit}(\#x{:},
      \langle\kw{identifier}(p)\rangle)),
      \kw{exprStmt}(\kw{identifier}(p))\rangle,\\
 \mathsf{setterExpansion}_{z}(x,R)&=
   \kw{message}_{z}\bigl(\kw{closure}_{z}
      (\langle\kw{formal}(p)\rangle,\epsilon,
       \mathsf{setterBody}_{z}(x,p)),\\[-1mm]
 &\hspace{38mm}\selsym{value:},\langle R\rangle\bigr).
\end{aligned}
\end{equation}
It therefore evaluates $R$ once, sends the setter, and returns the argument.

A cascade receiver is evaluated once.  For fresh synthetic closure identity
$q_c$, formal declaration $d_c$, and identifier $x_c$, define
\begin{equation}\label{eq:cascade-expansion}
\begin{aligned}
 &\mathsf{cascadeExpansion}_{z}
   (R,\bar c_0,
      \langle\kw{clause}(s_1,\bar R_1),\ldots,
             \kw{clause}(s_m,\bar R_m)\rangle)\\
 &\quad=\kw{message}_{z}\bigl(
   \kw{closure}_{q_c}(\langle\kw{formal}_{d_c}(x_c)\rangle,\epsilon,\\[-1mm]
 &\hspace{31mm}\langle
   \kw{exprStmt}(\mathsf{foldExplicit}_{z}
      (\kw{identifier}(x_c),\bar c_0))\rangle\cdot\\[-1mm]
 &\hspace{31mm}\langle
   \kw{exprStmt}(\kw{message}(\kw{identifier}(x_c),s_i,\bar R_i))
   \mid1\le i\le m\rangle),\\[-1mm]
 &\hspace{31mm}\selsym{value:},\langle R\rangle\bigr).
\end{aligned}
\end{equation}
The binding map inside the synthesized closure maps $x_c$ to $d_c$.

The remaining expression cases and statement elaboration are
preceded by normalization of local declarations:
\begin{equation}\label{eq:local-group-elaboration}
\begin{aligned}
 \mathsf{elabLocalDecl}_{\Gamma}
  (\kw{slot}_{d}(\alpha?,x,\kw{immutable},R))
   &=\kw{immutable}(d,\mathsf{elabExpr}_{\Gamma}(R)),\\
 \mathsf{elabLocalDecl}_{\Gamma}
  (\kw{slot}_{d}(\alpha?,x,\kw{mutable},R?))
   &=\kw{mutable}(d,\mathsf{elabOptionalExpr}_{\Gamma}(R?)),\\
 \mathsf{elabLocalGroup}_{\Gamma}(\kw{sequentialGroup}(\bar d))
   &=\langle\mathsf{elabLocalDecl}_{\Gamma}(d)\mid d\in\bar d\rangle,\\
 \mathsf{elabLocalGroup}_{\Gamma}(\kw{simultaneousGroup}(\bar d))
   &=\langle\kw{simultaneous}
      (\langle\mathsf{elabLocalDecl}_{\Gamma}(d)\mid d\in\bar d\rangle)\rangle,\\
 \mathsf{elabLocalGroups}_{\Gamma}(\epsilon)&=\epsilon,\\
 \mathsf{elabLocalGroups}_{\Gamma}(g\cdot\bar g)&=
   \mathsf{elabLocalGroup}_{\Gamma}(g)\cdot
   \mathsf{elabLocalGroups}_{\Gamma}(\bar g).
\end{aligned}
\end{equation}
\begin{equation}\label{eq:compound-expression-elaboration}
\begin{aligned}
 \Gamma_q&=\mathsf{activationContext}_{\Gamma}(q,\bar d,\bar g),\\
 \ell_q&=\mathsf{elabLocalGroups}_{\Gamma_q}(\bar g),\\
 b_q&=\mathsf{elabStatements}_{\Gamma_q}(\bar s),\\
 \mathsf{elabExpr}_{\Gamma}(\kw{closure}_{q}(\bar d,\bar g,\bar s))
   &=\kw{closureLiteral}(q,\mathsf{names}(\bar d),\ell_q,b_q),\\
 \mathsf{elabExpr}_{\Gamma}(\kw{object}_{L}(o_h))
   &=\kw{objectLiteral}(\mathsf{objectDescriptor}_{\Gamma}(L,o_h),
       \mathsf{objectSuperclass}_{\Gamma}(o_h)),\\
 \mathsf{elabExpr}_{\Gamma}(\kw{classExpr}_{D}(d))
   &=\mathsf{elabClass}_{\Gamma}(D,d),\\
 \mathsf{elabStatements}_{\Gamma}
   (\langle\kw{exprStmt}(R_i)\rangle_i)&=\\[-1mm]
 &\quad\langle\kw{expr}(\mathsf{elabExpr}_{\Gamma}(R_i))\rangle_i,\\
 \mathsf{elabStatements}_{\Gamma}
   (\bar s\cdot\langle\kw{returnStmt}(R)\rangle)
   &=\mathsf{elabStatements}_{\Gamma}(\bar s)\cdot\langle e_r\rangle,\\
 e_r&=\kw{return}(\mathsf{elabExpr}_{\Gamma}(R)).
\end{aligned}
\end{equation}
$\Gamma_q$ prepends the formal and local binder maps, sets the activation to
$q$, and enables a non-local return exactly when $q$ is a closure with an
enclosing method.  A return can only be the final surface statement by
Equation~\eqref{eq:newspeak-expression-peg}.

Access defaulting and the declaration cases are explicit:
\begin{equation}\label{eq:declaration-elaboration}
\begin{aligned}
 \mathsf{accessOf}(\kw{none})&=\Protected,\\
 \mathsf{accessOf}(\kw{public})&=\Public,\\
 \mathsf{accessOf}(\kw{protected})&=\Protected,\\
 \mathsf{accessOf}(\kw{private})&=\Private,\\
 \Gamma_f&=\mathsf{activationContext}_{\Gamma}(i_f,\bar d,\bar g),\\
 \ell_f&=\mathsf{elabLocalGroups}_{\Gamma_f}(\bar g),\\
 b_f&=\mathsf{elabStatements}_{\Gamma_f}(\bar s),\\
 a_f&=\mathsf{accessOf}(\alpha?),\\
 d_f&=\kw{method}_{i_f}(\alpha?,s,\bar d,\bar g,\bar s),\\
 d_s&=\kw{slot}_{d}(\alpha?,x,m,R?),\\
 d_l&=\kw{lazySlot}_{d}(\alpha?,x,m,R),\\
 d_n&=\kw{nestedClass}_{d}(\alpha?,D),\\
 \mathsf{elabDecl}_{\Gamma}(d_f)
   &=\langle i_f,s,a_f,\mathsf{names}(\bar d),
       \ell_f,b_f,\Gamma.K\rangle,\\
 \mathsf{elabDecl}_{\Gamma}(d_s)&=\\[-1mm]
 &\quad\mathsf{ordinarySlotDef}(d,x,a_f,m,
       \mathsf{elabOptionalExpr}_{\Gamma}(R?)),\\
 \mathsf{elabDecl}_{\Gamma}(d_l)&=\\[-1mm]
 &\quad\mathsf{lazySlotDef}(d,x,a_f,m,
       \mathsf{elabExpr}_{\Gamma}(R)),\\
 \mathsf{elabDecl}_{\Gamma}(d_n)&=\mathsf{nestedClassDef}(d,a_f,
       \mathsf{elabClass}_{\Gamma}(D,D)).
\end{aligned}
\end{equation}
$\Gamma_f$ prepends the method's formal and local declarations and enables
local return.  The ordinary- and lazy-slot constructors produce the physical
slot and the exact getter/setter bodies of
Section~\ref{sec:synthesized-methods}.  The nested-class constructor produces
its memoizing getter.  These constructors have different domains and names.

Class and object elaboration uses the message-template domain in
Equation~\eqref{eq:message-template-domain}.  Define
\begin{equation}\label{eq:front-end-message-template}
 \mathsf{elabMessageTemplate}_{\Gamma}
   (\kw{clause}(s,\bar R))=
   \langle s,\mathsf{elabExprList}_{\Gamma}(\bar R)\rangle.
\end{equation}
For a named class declaration $D$ with structure $h$, let $M_D$ and $M_D^c$
be the stable instance- and class-side mixins induced by $h$.  Its descriptor
constructor is
\begin{equation}\label{eq:front-end-body-descriptor}
\begin{aligned}
 \mathsf{bodyDescriptor}_{\Gamma}
   (D,x,u,s_f,\bar d,h,c)
   =\langle D,x,u,s_f,\mathsf{names}(\bar d),
      M_D,M_D^c,\mathsf{elabMessageTemplate}_{\Gamma}(c)\rangle.
\end{aligned}
\end{equation}
This is exactly the tuple $\chi_b$ of Equation~\eqref{eq:body-descriptor}.
Omitted inheritance is
\begin{equation}\label{eq:front-end-default-superclass}
\begin{aligned}
 \mathsf{defaultSuperclass}_{\Gamma}&=
  \begin{cases}
   \Object&\Gamma.K=\kw{top},\\
   \kw{implicit}(\#\kw{Object},\epsilon,
      \mathsf{bindingAnnotation}_{\Gamma}(\kw{Object}))&\text{otherwise,}
  \end{cases}\\
 \mathsf{defaultInitializer}&=\kw{clause}(\#\kw{new},\epsilon).
\end{aligned}
\end{equation}
The two non-application class cases are therefore
\begin{equation}\label{eq:front-end-class-body-elaboration}
\begin{aligned}
 &\mathsf{elabClass}_{\Gamma}
  (D,\kw{classDecl}(x,s_f,\bar d,\kw{defaultInheritance}(h)))\\
 &\quad=\kw{classBody}(\chi_b,\mathsf{defaultSuperclass}_{\Gamma}),\\[-1mm]
 &\qquad\chi_b=\mathsf{bodyDescriptor}_{\Gamma}
      (D,x,\kw{implicitSuperclass},s_f,\bar d,h,
       \mathsf{defaultInitializer}),\\
 &\mathsf{elabClass}_{\Gamma}
  (D,\kw{classDecl}(x,s_f,\bar d,
       \kw{explicitInheritance}(R,c,h)))\\
 &\quad=\kw{classBody}
   (\chi_b,\mathsf{elabExpr}_{\Gamma}(R)),\\[-1mm]
 &\qquad\chi_b=\mathsf{bodyDescriptor}_{\Gamma}
      (D,x,\kw{explicitSuperclass},s_f,\bar d,h,c).
\end{aligned}
\end{equation}

For a mixin chain, write its first inheritance head as
$h_0=\kw{inheritanceHead}(R_0,c_0)$ and its nonempty suffix as
$\bar h=\langle h_1,\ldots,h_k\rangle$, where
$h_j=\kw{inheritanceHead}(R_j,c_j)$.  Each application uses $c_{j-1}$ to
initialize its superclass and $c_j$ to initialize the newly applied mixin.
The synthetic identities and names required for an intermediate application
are deterministic injections of the enclosing declaration identity and its
one-based position:
\begin{equation}\label{eq:front-end-application-identities}
\begin{aligned}
 D_{A,j}&=\mathsf{applicationDeclId}(D,j),&
 i_{A,j}&=\mathsf{applicationInitId}(D,j),\\
 n_{A,j}&=\mathsf{applicationName}(D,j),&
 \bar x_{A,j}&=\mathsf{applicationFormals}(D,j,\selector(c_j)).
\end{aligned}
\end{equation}
The names in $\bar x_{A,j}$ are pairwise fresh, and their number is the arity
of $\selector(c_j)$.  Let $M^c_{A,j}$ be the stable class-side mixin containing
only the synthesized public primary factory for that signature.  The complete
descriptor of an intermediate application is
\begin{equation}\label{eq:front-end-anonymous-application-descriptor}
\begin{aligned}
 \mathsf{anonAppDesc}_{\Gamma}(D,j,c_{j-1},c_j)
   &=\langle D_{A,j},n_{A,j},\selector(c_j),\bar x_{A,j},\\[-1mm]
   &\quad M^c_{A,j},i_{A,j},m_{j-1},m_j\rangle,\\
 m_{j-1}&=\mathsf{elabMessageTemplate}_{\Gamma}(c_{j-1}),\\
 m_j&=\mathsf{elabMessageTemplate}_{\Gamma}(c_j).
\end{aligned}
\end{equation}
For the last application of a body-less named class, the declaration's own
identity, name, factory formals, and class-side mixin replace those synthetic
components:
\begin{equation}\label{eq:front-end-named-application-descriptor}
\begin{aligned}
 &\mathsf{namedAppDesc}_{\Gamma}
   (D,x,s_f,\bar d,c_{k-1},c_k)={}\\
 &\quad\langle D,x,s_f,\mathsf{names}(\bar d),M_D^c,i_{A,k},
   \mathsf{elabMessageTemplate}_{\Gamma}(c_{k-1}),
   \mathsf{elabMessageTemplate}_{\Gamma}(c_k)\rangle.
\end{aligned}
\end{equation}
The anonymous prefix fold returns both the resulting class expression and the
last initializer clause; retaining the latter makes the next superclass
initializer explicit:
\begin{equation}\label{eq:front-end-mixin-fold}
\begin{aligned}
 \mathsf{foldAnonymous}_{\Gamma,D}(E,c_0,\epsilon)&=\langle E,c_0\rangle,\\
 \mathsf{foldAnonymous}_{\Gamma,D}
  (E,c_{j-1},\langle h_j\rangle\cdot\bar h)
   &=\mathsf{foldAnonymous}_{\Gamma,D}(E_j,c_j,\bar h),\\
 E_j&=\kw{mixinApply}(\chi_{A,j},E,
       \mathsf{elabExpr}_{\Gamma}(R_j)),\\
 \chi_{A,j}&=\mathsf{anonAppDesc}_{\Gamma}(D,j,c_{j-1},c_j).
\end{aligned}
\end{equation}
A body-less chain reserves its last application for the named descriptor.  Put
\begin{equation}\label{eq:front-end-bodyless-prefix}
 \mathsf{foldAnonymous}_{\Gamma,D}
  (\mathsf{elabExpr}_{\Gamma}(R_0),c_0,
   \langle h_1,\ldots,h_{k-1}\rangle)
   =\langle E_{k-1},c_{k-1}\rangle.
\end{equation}
Then
\begin{equation}\label{eq:front-end-bodyless-mixin-elaboration}
\begin{aligned}
 &\mathsf{elabClass}_{\Gamma}
  (D,\kw{classDecl}(x,s_f,\bar d,
     \kw{mixinChain}(R_0,c_0,\bar h,\kw{none})))={}\\
 &\quad\kw{mixinApply}(
   \mathsf{namedAppDesc}_{\Gamma}
      (D,x,s_f,\bar d,c_{k-1},c_k),
   E_{k-1},\mathsf{elabExpr}_{\Gamma}(R_k)).
\end{aligned}
\end{equation}
This equation is defined only when $s_f=\selector(c_k)$ and $k\ge1$, as the
grammar requires.  If the chain ends in a body, put
\begin{equation}\label{eq:front-end-bodied-prefix}
 \mathsf{foldAnonymous}_{\Gamma,D}
  (\mathsf{elabExpr}_{\Gamma}(R_0),c_0,\bar h)=\langle E_k,c_k\rangle.
\end{equation}
Then
\begin{equation}\label{eq:front-end-bodied-mixin-elaboration}
\begin{aligned}
 &\mathsf{elabClass}_{\Gamma}
  (D,\kw{classDecl}(x,s_f,\bar d,
     \kw{mixinChain}(R_0,c_0,\bar h,h)))={}\\
 &\quad\kw{classBody}
   (\mathsf{bodyDescriptor}_{\Gamma}
      (D,x,\kw{explicitSuperclass},s_f,\bar d,h,c_k),E_k).
\end{aligned}
\end{equation}
Equations~\eqref{eq:front-end-anonymous-application-descriptor} and
\eqref{eq:front-end-named-application-descriptor} are instances of the complete
application tuple $\chi_a$ in Equation~\eqref{eq:application-descriptor}; no
descriptor component is supplied informally.

An object literal uses the same body derivation with an object-literal identity
$L$, the public nullary synthetic class-side factory, and the optional header's
inheritance receiver and initializer clause:
\begin{equation}\label{eq:front-end-object-elaboration}
\begin{aligned}
 m_0&=\mathsf{elabMessageTemplate}_{\Gamma}
         (\mathsf{defaultInitializer}),\\
 m_1&=\mathsf{elabMessageTemplate}_{\Gamma}(c),\\
 \mathsf{objectDescriptor}_{\Gamma}
  (L,\kw{implicitObjectHeader}(h))
   &=\langle L,\kw{implicitSuperclass},M_{L,h},M^c_{L,h},m_0\rangle,\\
 \mathsf{objectDescriptor}_{\Gamma}
  (L,\kw{explicitObjectHeader}(R,c,h))
   &=\langle L,\kw{explicitSuperclass},M_{L,h},M^c_{L,h},m_1\rangle,\\
 \mathsf{objectSuperclass}_{\Gamma}(\kw{implicitObjectHeader}(h))
   &=\mathsf{defaultSuperclass}_{\Gamma},\\
 \mathsf{objectSuperclass}_{\Gamma}(\kw{explicitObjectHeader}(R,c,h))
   &=\mathsf{elabExpr}_{\Gamma}(R).
\end{aligned}
\end{equation}
Equation~\eqref{eq:front-end-object-elaboration} is exactly the descriptor
$\chi_o$ in Equation~\eqref{eq:object-literal-descriptor}.

For an identified unit $S$, the maps consumed by $P$ are now defined directly.
Let $\mathsf{ancestors}_{S}(n)$ be the strict lexical-ancestor sequence of
node $n$, nearest first.  Let $M_K$ be the stable mixin identity derived from
class-body identity $K$.  Then
\begin{equation}\label{eq:derived-program-indexes}
\begin{aligned}
 \mathsf{sourceMap}_{S}(d)&=\mathsf{subtree}_{S}(d),\\
 \mathsf{declMixinMap}_{S}(K)&=M_K,\\
 \mathsf{lexParentMap}_{S}(D)&=
   \mathsf{firstClass}(\mathsf{ancestors}_{S}(D)),\\
 \mathsf{bodyMap}_{S}(i_f)&=
   \mathsf{elabStatements}_{\Gamma_{i_f}}(\mathsf{surfaceBody}_{S}(i_f)),\\
 \mathsf{scopeChainMap}_{S}(z)&=
   \langle d\in\mathsf{ancestors}_{S}(z)\mid
      d\in\mathit{ClassBodyDeclId}\uplus\mathit{ActDeclId}\rangle,\\
 \mathsf{activationMemberMap}_{S}(q)&=
   \{\#x\mid x\text{ is a formal or local declared directly by }q\},\\
 \mathsf{ownerMap}_{S}(K)&=
  \begin{cases}
   \kw{activationOwner}&\mathsf{firstOwner}_{S}(K)\in\mathit{ActDeclId},\\
   \kw{classOwner}&\mathsf{firstOwner}_{S}(K)\in\mathit{ClassDeclId},\\
   \kw{objectLiteralOwner}&\mathsf{firstOwner}_{S}(K)\in\mathit{ObjLitDeclId},\\
   \kw{topOwner}&\mathsf{firstOwner}_{S}(K)=\kw{none}.
  \end{cases}
\end{aligned}
\end{equation}
The mixin entry itself is
\begin{equation}\label{eq:derived-mixin-entry}
\begin{aligned}
 \phi_K&=\mathsf{explicitMethods}_{S}(K)
          \uplus\mathsf{synthMethods}_{S}(K),\\
 \delta_K&=\mathsf{physicalSlots}_{S}(K),
 &\nu_K&=\mathsf{nestedDecls}_{S}(K),\\
 \iota_K&=\mathsf{initializer}_{S}(K),
 &\mathsf{mixinMap}_{S}(M_K)&=
   \langle K,\phi_K,\delta_K,\nu_K,\iota_K\rangle.
\end{aligned}
\end{equation}
The disjoint union is defined only when selectors are unique.
$\mathsf{initializer}_{S}(K)$ is the tuple in
Equation~\eqref{eq:mixin-initializer}: its factory selector and formals come
from the class header, its source-ordered groups are transformed by
Equation~\eqref{eq:declaration-elaboration}, and its body by
Equation~\eqref{eq:compound-expression-elaboration}.

Top-level actor seeds and the one underspecified pattern form are
\begin{equation}\label{eq:derived-auxiliary-indexes}
\begin{aligned}
 \mathsf{actorSeedMap}_{S}(M_K)&=
   \langle K,M_K,\mathsf{factoryDescriptor}_{S}(K)\rangle
   &&\mathsf{ownerMap}_{S}(K)=\kw{topOwner},\\
 \mathsf{patternVariableMap}_{S}(z)&=
   \mathsf{patternVariableParameter}_{S}(z)
   &&\kw{variablePat}_{z}(x)\in S.
\end{aligned}
\end{equation}
The latter is precisely the explicit parameter isolated by
Equation~\eqref{eq:variable-pattern-parameter}; it is not used for any other
surface form.

Collecting Equations~\eqref{eq:derived-program-indexes}--
\eqref{eq:derived-auxiliary-indexes}, define
\begin{equation}\label{eq:declaration-derivation}
\begin{aligned}
 \mathsf{deriveProgram}(S)=\langle P,e\rangle
 \Longleftrightarrow{}&
 &P.\mathit{unit}=S\land\\
 &P.\mathit{source}=\mathsf{sourceMap}_{S}\land\\
 &P.\mathit{declMixin}=\mathsf{declMixinMap}_{S}\land\\
 &P.\mathit{mixin}=\mathsf{mixinMap}_{S}\land\\
 &P.\mathit{lexParent}=\mathsf{lexParentMap}_{S}\land\\
 &P.\mathit{body}=\mathsf{bodyMap}_{S}\land\\
 &P.\mathit{scopeChain}=\mathsf{scopeChainMap}_{S}\land\\
 &P.\mathit{activationMembers}=\mathsf{activationMemberMap}_{S}\land\\
 &P.\mathit{owner}=\mathsf{ownerMap}_{S}\land\\
 &P.\mathit{actorSeed}=\mathsf{actorSeedMap}_{S}\land\\
 &P.\mathit{patternVariable}=\mathsf{patternVariableMap}_{S}\land\\
 &e=\mathsf{elabTop}_{S}(S).
\end{aligned}
\end{equation}
The unchanged atom and platform maps are supplied as parameters to
$\mathsf{deriveProgram}$ when a program is installed.

The static predicates are defined without an informal residual clause:
\begin{equation}\label{eq:static-well-formedness-components}
\begin{aligned}
 \mathsf{bindingsOK}(S)&\Longleftrightarrow
   \bigl(\forall\Gamma,B_i\in\Gamma.\bar B.\;
      \mathsf{declaredNames}(B_i)\text{ are pairwise distinct}\bigr)\land\\[-1mm]
 &\quad\bigl(\forall\kw{outer}_{z}(x)\in S.\;
      \mathsf{classBind}_{\Gamma_z}(x)\ne\undef\land
      \Gamma_z.D_0\ne\kw{none}\bigr)\land\\[-1mm]
 &\quad\bigl(\forall\kw{self}_{z}\text{ used as receiver in }S.\;
      \Gamma_z.D_0\ne\kw{none}\bigr),\\
 \mathsf{membersOK}(S)&\Longleftrightarrow
   \forall K,s.\;|\{f\mid f\text{ is induced directly in }K\land
                         \selector(f)=s\}|\le1,\\
 \mathsf{arityOK}(S)&\Longleftrightarrow
   \forall f.\;|f.\bar x|=\mathsf{arity}(\selector(f)),\\
 \mathsf{returnsOK}(S)&\Longleftrightarrow
   \forall R\in\mathit{ReturnNode}(S).\;
   (\mathsf{local}(R)\Rightarrow r_t(R))\land
   (\mathsf{nonlocal}(R)\Rightarrow h_t(R)),\\
 \mathsf{contextsOK}(S)&\Longleftrightarrow
   \mathsf{bindingsOK}(S)\land\mathsf{returnsOK}(S)\land
   \forall R\in\mathit{SuperSendNode}(S).\;q(R)\ne\kw{none},\\
 \mathsf{surfaceOK}(S)&\Longleftrightarrow
   \mathsf{idsOK}(S)\land\mathsf{membersOK}(S)\land
   \mathsf{arityOK}(S)\land\mathsf{contextsOK}(S)\land
   \mathsf{shapesOK}(S).
\end{aligned}
\end{equation}
The surface-shape condition is
\begin{equation}\label{eq:surface-context-validity}
\begin{split}
 \mathsf{lastClause}
  (\langle h_1,\ldots,h_k\rangle)&=c_k
   \quad\text{when }h_k=\kw{inheritanceHead}(R_k,c_k),\\
 \mathsf{shapesOK}(S)\Longleftrightarrow{}&
  \text{every }\kw{outer}_{z}(x)\text{ or }\kw{super}_{z}
     \text{ is the receiver of a message}\ \land\\
 &\text{every return is final in its statement sequence}\ \land\\
 &\text{every immutable slot has an initializer}\ \land\\
 &\text{every lazy slot has an initializer}\ \land\\
 &\forall\kw{classDecl}(x,s_f,\bar d,
       \kw{mixinChain}(R_0,c_0,\bar h,\kw{none}))\in S.\;
       s_f=\selector(\mathsf{lastClause}(\bar h)).
\end{split}
\end{equation}
Ordinary unbound implicit sends are permitted by
Equation~\eqref{eq:implicit-binding-annotation}; named outer targets and
duplicate declarations in one binder map are rejected by the first predicate.

The source-to-core judgment is
\begin{equation}\label{eq:source-to-annotated-core}
 \begin{gathered}
 G,\kappa,u\vdash\mathsf{frontEnd}\Downarrow
   \langle P,e,\kappa'\rangle\\[-1mm]
 \Longleftrightarrow\\[-1mm]
  \mathsf{parse}_{G}(u)=t,\quad
  \mathsf{ast}_{G,u}(t)=\langle R_0,\bar m\rangle,\quad
  R=\mathsf{attachMetadata}_{u}(R_0,\bar m),\quad
  \mathsf{identify}_{\kappa}(R)=\langle S,\kappa'\rangle,\\
  \mathsf{surfaceOK}(S),\quad
  \mathsf{deriveProgram}(S)=\langle P,e\rangle,\quad
  \mathsf{TL}_{P}(e).
 \end{gathered}
\end{equation}
Core validity must cover expressions stored throughout $P$, not only the
top-level expression.  Let $\mathsf{coreRoots}(P,e)$ be the least finite set
containing $e$ and every expression occurring in method and closure bodies,
their local initializers, mixin initializer groups and bodies, pattern-variable
expansions, and the deferred superclass or mixin messages stored in class,
mixin-application, object-literal, and actor-seed descriptors:
\begin{equation}\label{eq:annotated-core-roots}
\begin{split}
 \mathsf{coreRoots}(P,e)=\mathsf{sub}\bigl(&\{e\}
   \cup P.\mathit{methodBodies}\cup P.\mathit{methodLocals}\\
 &\cup P.\mathit{closureBodies}\cup P.\mathit{closureLocals}\\
 &\cup P.\mathit{mixinInitializers}\cup P.\mathit{patternExpansions}\\
 &\cup P.\mathit{classMessages}\cup P.\mathit{mixinMessages}\\
 &\cup P.\mathit{objectMessages}\cup P.\mathit{actorMessages}\bigr).
\end{split}
\end{equation}
Here $\mathsf{sub}$ recursively includes tuple, pattern, cascade, send, lazy,
class, mixin-application, object-literal, and nested-class subexpressions.  For
Equation~\eqref{eq:annotated-core-roots}, $P.\mathit{mixinInitializers}$
projects the groups and bodies of all mixin initializers,
$P.\mathit{patternExpansions}$ is the range of the finite pattern-variable
table, and the four $\mathit{Messages}$ projections collect the message
templates in the corresponding finite descriptor tables.  These are
abbreviations for existing Program fields, not new semantic stores.  For
fuel $k$, $\mathsf{annOK}^{k}_{P}(r,b)$ checks every declaration and scope
annotation in $r$, where $b$ says whether a super send is permitted.  In
particular, it checks that named class and closure declarations occur in the
appropriate finite Program table, that an annotated implicit send is declared
by its recorded scope, and that a super send occurs only in installed class
code.  The executable gate is
\begin{equation}\label{eq:annotated-core-checker}
\begin{split}
 \mathsf{coreCheck}_{k}(P,e)=\kw{true}\Longleftrightarrow{}&
 P.\mathit{declMixin}\text{ is injective}\ \land\\
 &\forall r\in\mathsf{coreRoots}(P,e)\setminus\{e\}.\;
       \mathsf{annOK}^{k}_{P}(r,\kw{true})\ \land\\
 &\mathsf{annOK}^{k}_{P}(e,\kw{false})\ \land\ \mathsf{TL}_{P}(e).
\end{split}
\end{equation}
Core validity, used by loaders and reflective validation, is therefore the
independently checkable predicate
\begin{equation}\label{eq:annotated-core-validity}
\begin{split}
 \mathsf{coreOK}(P,e)\Longleftrightarrow{}&
 \exists k>0.\ \mathsf{coreCheck}_{k}(P,e)=\kw{true}\ \land\\
 &\mathsf{deriveProgram}(P.\mathit{unit})=\langle P,e\rangle\ \land
   \mathsf{surfaceOK}(P.\mathit{unit}).
\end{split}
\end{equation}
For reflective validation, which replaces a program but does not itself start
its top-level expression, define the projection
\begin{equation}\label{eq:annotated-program-validity}
 \mathsf{coreProgramOK}(P)
 \Longleftrightarrow
 \exists e.\;\mathsf{deriveProgram}(P.\mathit{unit})=\langle P,e\rangle
              \land\mathsf{coreOK}(P,e).
\end{equation}
Equations~\eqref{eq:source-to-annotated-core} and
\eqref{eq:annotated-core-validity} are the promised static connection from the
Specification's PEG AST to every declaration and send annotation consumed by
the dynamic semantics.  Moreover, determinism of PEG recognition and the
functional equations above give the front-end uniqueness property
\begin{equation}\label{eq:front-end-uniqueness}
 \frac{G_{\NS},\kappa,u\vdash\mathsf{frontEnd}\Downarrow X\qquad
       G_{\NS},\kappa,u\vdash\mathsf{frontEnd}\Downarrow Y}{X=Y}.
\end{equation}
Totality holds on exactly the inputs that parse completely and satisfy
$\mathsf{surfaceOK}$; failure at any earlier partial function is a static
error, never a stuck core term.

\subsection{Reified language services}
\label{sec:reified-front-ends}

The host registry already introduced in Equation~\eqref{eq:compilation-unit}
now has the concrete type
\begin{equation}\label{eq:language-registry}
 \mathcal L:\mathit{LanguageId}\rightharpoonup\mathit{ObjRef}.
\end{equation}
$\mathcal L(\lambda)$ is an ordinary Newspeak language object.  It answers the
public message \texttt{frontEndImage} with an immutable value
\begin{equation}\label{eq:front-end-image}
 I_L=\kw{frontEndImage}\langle o_p,o_e,o_c,o_s,G,v_L\rangle,
\end{equation}
whose components are parser, elaborator, compiler, and Simulator objects, the
declarative PEG value, and a language-defined revision value.  This record is
a protocol value, not a new privileged heap-record kind.  Its component
objects have ordinary classes, methods, enclosing objects, and slots.

For a caller actor $A$, write
\begin{equation}\label{eq:ordinary-service-call}
 \Omega,A\vdash\mathsf{call}(o,s,\bar v)
    \Downarrow\langle\Omega',v\rangle
\end{equation}
iff the ordinary public send of $\langle s,\bar v\rangle$ to $o$ starts in
$A$, follows the actor and sequential transition relations already defined,
and returns $v$ to its caller.  A signaled exception or broken reply produces
the corresponding failure rather than a successful judgment.  Thus
Equation~\eqref{eq:ordinary-service-call} is notation for the reflexive
transitive closure of existing rules, not a second evaluator.

Surface ASTs and annotated-core images are ordinary immutable Newspeak value
graphs.  Executable artifacts are immutable values as well.  The VM boundary
supplies injective partial
decoders
\begin{equation}\label{eq:front-end-value-decoders}
\begin{aligned}
 \mathsf{decodeSurface}&:\mathit{ObjRef}\rightharpoonup\mathit{SurfaceAST},\\
 \mathsf{decodeCore}&:\mathit{ObjRef}\rightharpoonup
        (\mathit{Program}\times\mathit{CoreExpression}),\\
 \mathsf{decodeArtifact}&:\mathit{ObjRef}\rightharpoonup\mathit{Artifact}.
\end{aligned}
\end{equation}
They inspect representation and allocate no object, invoke no user method, and
accept only finite immutable graphs.  Their inverses are ordinary constructors
provided by the language implementation.

Admission is attached to the values obtained at a use boundary, not to the
permanent identity of a component object.  A front-end image is admitted when
all four component references are live ordinary Newspeak objects, its grammar
passes Equation~\eqref{eq:peg-admissibility-certificate}, and successful
elaboration satisfies the complete checker:
\begin{equation}\label{eq:front-end-artifact-admission}
\begin{split}
 \mathsf{frontEndAdmitted}_{\Omega}(I_L)\Longleftrightarrow{}&
 I_L=\kw{frontEndImage}\langle o_p,o_e,o_c,o_s,G,v_L\rangle\ \land\\
 &\bigwedge_{o\in\{o_p,o_e,o_c,o_s\}}\mathsf{liveOrdinary}_{\Omega}(o)
 \ \land\ \mathsf{pegCertOK}(G)=\kw{true}\ \land\\
 &\forall S,P,e.\;
   \mathsf{surfaceOK}(S)\land
   \mathsf{deriveProgram}(S)=\langle P,e\rangle\\
 &\hspace{24mm}\Rightarrow\exists k>0.\;
       \mathsf{coreCheck}_{k}(P,e)=\kw{true}.
\end{split}
\end{equation}
Here $\mathsf{liveOrdinary}_{\Omega}(o)$ means that $o$ is present in the
ordinary-object component of the appropriate heap in $\Omega$.
The compiler's product is separately admitted by
Equation~\eqref{eq:artifact-validity}; the Simulator's code/PC representation
is separately admitted by the certificate in
Equation~\eqref{eq:simulator-step-correspondence}.  Replacing a component
therefore grants no authority to bypass the fixed core boundary.

Capture one front-end image and run its three stages by ordinary sends:
\begin{equation}\label{eq:metacircular-compilation-pipeline}
\begin{split}
 &\Omega_0,A,\lambda,u\vdash\mathsf{compileUnit}
      \Downarrow\langle\Omega_4,B,P,e,I_L\rangle
 \quad\Longleftrightarrow\\
 &\quad o_L=\mathcal L(\lambda)\ \land\
 \Omega_0,A\vdash\mathsf{call}(o_L,\selsym{frontEndImage},\epsilon)
      \Downarrow\langle\Omega_1,I_L\rangle\ \land\\
 &\quad I_L=\kw{frontEndImage}\langle o_p,o_e,o_c,o_s,G,v_L\rangle\ \land\\
 &\quad \mathsf{frontEndAdmitted}_{\Omega_1}(I_L)\ \land\
 \Omega_1,A\vdash\mathsf{call}(o_p,\selsym{parse:},\langle u\rangle)
      \Downarrow\langle\Omega_2,v_s\rangle\ \land\\
 &\quad \mathsf{decodeSurface}(v_s)=S\ \land\mathsf{idsOK}(S)\ \land\
 \Omega_2,A\vdash\mathsf{call}(o_e,\selsym{elaborate:},\langle v_s\rangle)
      \Downarrow\langle\Omega_3,v_k\rangle\ \land\\
 &\quad \mathsf{decodeCore}(v_k)=\langle P,e\rangle\ \land
       \mathsf{coreOK}(P,e)\ \land\
 \Omega_3,A\vdash\mathsf{call}(o_c,\selsym{compile:},\langle v_k\rangle)
      \Downarrow\langle\Omega_4,v_b\rangle\ \land\\
 &\quad \mathsf{decodeArtifact}(v_b)=B\ \land
       \mathsf{artifactOK}(B,P,e).
\end{split}
\end{equation}
The source string $u$ is represented by the platform's immutable String value;
the notation passes its contents rather than granting the compiler a raw host
pointer.  A language implementation may combine stages internally, but its
observable result must refine Equation~\eqref{eq:metacircular-compilation-pipeline}.

An artifact contains the core image, concrete code images, relocation data,
source/control maps, and the captured language revision:
\begin{equation}\label{eq:artifact-domain}
 B=\langle P,e,\mathcal B,\mathsf{reloc},\mathsf{controlMap},v_L\rangle.
\end{equation}
Its validity is
\begin{equation}\label{eq:artifact-validity}
\begin{split}
 \mathsf{artifactOK}(B,P,e)\Longleftrightarrow{}&
 \mathsf{coreOK}(P,e)\land B.P=P\land B.e=e\land\\
 &\dom(B.\mathsf{controlMap})=\mathsf{pcs}(B.\mathcal B)\land\\
 &\mathsf{targetsOK}(B.\mathcal B,B.\mathsf{reloc})\land
 \mathsf{codeRefinesCore}(B.\mathcal B,P,e).
\end{split}
\end{equation}
The final conjunct of Equation~\eqref{eq:artifact-validity} is the
compiler-correctness obligation.  Concrete execution, observed at the
source/control map, must weakly simulate the sequential and actor semantics.
A changed compiler may accept new surface forms or omit old conveniences by
changing parsing and elaboration, but an artifact admitted by the VM must
still satisfy the fixed core contract.  Changing the instruction set or
primitive meanings requires a VM refinement, not merely a language front-end
change.

The earlier compilation-unit judgment is therefore the projection
\begin{equation}\label{eq:compilation-unit-projection}
 \mathcal L,\Omega,A\vdash\langle\lambda,u\rangle
      \mathsf{compilesTo}\langle P,e,B\rangle
 \quad\Longleftrightarrow\quad
 \exists\Omega',I_L.\;
 \Omega,A,\lambda,u\vdash\mathsf{compileUnit}
      \Downarrow\langle\Omega',B,P,e,I_L\rangle.
\end{equation}

\subsection{Ordinary change, phase behavior, and bootstrap}
\label{sec:metacircular-change}

Let $\mathsf{components}_{\Omega}(\lambda)$ be the value obtained by the first
call in Equation~\eqref{eq:metacircular-compilation-pipeline}.  A successful
existing reflective commit may change methods of $o_L,o_p,o_e,o_c,o_s$, or may
change ordinary slots from which \texttt{frontEndImage} constructs its result.
Its effect on the next compilation is simply
\begin{equation}\label{eq:front-end-reflective-change}
 \Omega\xrightarrow{\kw{commitReflection}(A,\Delta,n+1)}\Omega'
 \ \land\
 \mathsf{components}_{\Omega}(\lambda)\ne
 \mathsf{components}_{\Omega'}(\lambda).
\end{equation}
Equation~\eqref{eq:front-end-reflective-change} is a consequence that an
ordinary reflection transition may have; it is not a new transition rule and
there is no $\mathsf{installParser}$ or $\mathsf{installCompiler}$ command.
Likewise, assigning a new component object through the language object's
ordinary mutable protocol needs no reflection at all when the caller already
has that authority.

The phase behavior follows from existing activation and dispatch rules:
\begin{equation}\label{eq:front-end-phase-behavior}
\begin{aligned}
 \mathsf{capturedImage}(c)&=I_L
   &&\text{for a compilation }c\text{ already past its first call},\\
 \mathsf{currentTerm}_{\Omega'}(a)&=\mathsf{currentTerm}_{\Omega}(a)
   &&\text{for every materialized stage activation }a,\\
 \mathsf{dispatchProgram}(a')&=P'
   &&\text{for a later send dispatched after the commit.}
\end{aligned}
\end{equation}
Thus self-replacement needs no circular premise.  Source defining a replacement
parser is parsed and compiled by the currently executing language objects; an
ordinary commit changes their code or references; subsequent compilation sees
the replacement.  A replacement not expressible in the old accepted language
must enter as an artifact already satisfying Equation~\eqref{eq:artifact-validity}.

Finally, metacircularity has a finite base.  An initial image supplies
\begin{equation}\label{eq:bootstrap-basis}
 \mathsf{bootstrap}=\langle
   \Omega_0,\mathcal L_0,
   \mathsf{decodeSurface},\mathsf{decodeCore},\mathsf{decodeArtifact},
   \mathsf{loadCode},\mathsf{primitiveStep}\rangle.
\end{equation}
After bootstrap, every object reachable through $\mathcal L_0$ is an ordinary
Newspeak object and can be changed by the mechanisms above.  The decoders,
artifact loader, concrete instruction dispatch, allocation, and host I/O remain
VM parameters unless a particular implementation supplies a Newspeak
refinement of them.  This boundary states what is not metacircular without
preventing a more completely bootstrapped VM from discharging those parameters.
